% PDF page 232
\source{gilbarg-trudinger:page-0232}

\pagestyle{empty}


\section*{9.1. Maximum Principles for Strong Solutions}

In this section, we treat the extension of the classical maximum principle in Chapter 3 to strong solutions, in particular to solutions in the Sobolev space $W^{2,n}_{\mathrm{loc}}(\Omega)$. Recall that an operator $L$ of the form (9.1) is elliptic in the domain $\Omega$ if the coefficient matrix $\mathcal{A}=[a^{ij}]$ is positive everywhere in $\Omega$. For such operators we will let $\mathcal{D}$ denote the determinant of $\mathcal{A}$ and set $\mathcal{D}^{*}=\mathcal{D}^{1/n}$ so that $\mathcal{D}^{*}$ is the geometric mean of the eigenvalues of $\mathcal{A}$ and

\[
0<\lambda\le \mathcal{D}^{*}\le \Lambda
\]

where as before $\lambda$, $\Lambda$ denote, respectively, the minimum and maximum eigenvalues of $\mathcal{A}$. Our conditions on the coefficients of $L$ and inhomogeneous term $f$ in the equation (9.2) will now take the form

\[
\text{(9.3)}\qquad |b|/\mathcal{D}^{*},\, f/\mathcal{D}^{*}\in L^{n}(\Omega),\qquad c\le 0 \quad \text{in }\Omega.
\]

The following weak maximum principle, of A.\ D.\ Aleksandrov, can now be formulated as an extension of the apriori bound, Theorem 3.7.

\textbf{Theorem 9.1.}
Let $Lu\ge f$ in a bounded domain $\Omega$ and $u\in C^{0}(\overline{\Omega})\cap W^{2,n}_{\mathrm{loc}}(\Omega)$. Then

\[
\text{(9.4)}\qquad \sup_{\Omega}u\le \sup_{\partial\Omega}u^{+}+C\|f/\mathcal{D}^{*}\|_{L^{n}(\Omega)}
\]

where $C$ is a constant depending only on $n$, diam $\Omega$ and $\|b/\mathcal{D}^{*}\|_{L^{n}(\Omega)}$.

The Sobolev embedding theorem, in particular Corollary 7.11, guarantees that functions in $W^{2,n}_{\mathrm{loc}}(\Omega)$ are at least continuous in $\Omega$. If $u$ is not also assumed continuous on $\overline{\Omega}$ in the hypotheses of Theorem 9.1, the conclusion (9.4) can be modified by replacing $\sup_{\partial\Omega}u^{+}$ by $\limsup_{x\to\partial\Omega}u^{+}$.

The proof of Theorem 9.1 depends upon the notions of contact set and normal mapping and certain aspects of it will be important for later considerations. If $u$ is an arbitrary continuous function on $\Omega$, we define the upper contact set of $u$, denoted $\Gamma^{+}$ or $\Gamma^{+}_{u}$, to be the subset of $\Omega$ where the graph of $u$ lies below a support hyperplane in $\mathbb{R}^{n+1}$, that is,

\[
\text{(9.5)}\qquad \Gamma^{+}=\{y\in\Omega\mid u(x)\le u(y)+p\cdot(x-y)\ \text{for all }x\in\Omega,\ \text{for some }p=p(y)\in\mathbb{R}^{n}\}.
\]

Clearly, $u$ is a concave function in $\Omega$ if and only if $\Gamma^{+}=\Omega$. When $u\in C^{1}(\Omega)$ we must have $p=Du(y)$ in (9.5) as any support hyperplane must then be a tangent hyperplane to the graph of $u$. Furthermore, when $u\in C^{2}(\Omega)$, the Hessian matrix $D^{2}u=[D_{ij}u]$ is nonpositive on $\Gamma^{+}$. In general, the set $\Gamma^{+}$ is closed relative to $\Omega$.


% PDF page 233
\source{gilbarg-trudinger:page-0233}

For an arbitrary function $u \in C^0(\Omega)$ we define the \emph{normal mapping}, $\chi(y) = \chi_u(y)$ of a point $y \in \Omega$ to be the set of ``slopes'' of support hyperplanes at $y$ lying above the graph of $u$, that is,

\begin{equation}
\chi(y) = \{p \in \R^n \mid u(x) \leq u(y) + p \cdot (x - y) \quad \text{for all } x \in \Omega\}.
\tag{9.6}
\end{equation}

Clearly $\chi(y)$ is nonempty if and only if $y \in \Gamma^+$. Furthermore, when $u \in C^1(\Omega)$, then $\chi(y) = Du(y)$ on $\Gamma^+$; that is, $\chi$ is the gradient vector field of $u$ on $\Gamma^+$. As a useful example of a nondifferentiable function $u$, let us take $\Omega$ to be a ball $B = B_R(z)$ and $u$ to be the function whose graph is a cone with base $\Omega$ and vertex $(z, a)$ for some positive $a \in \R$, that is,

\[
u(x) = a\left(1 - \frac{\abs{x - z}}{R}\right).
\]

Then, we have

\begin{equation}
\chi(y) =
\begin{cases}
\dfrac{-a(y - z)}{R \abs{y - z}} & \text{for } y \neq z \\
B_{a/R}(0) & \text{for } y = z.
\end{cases}
\tag{9.7}
\end{equation}

First, we prove:

\medskip

\textbf{Lemma 9.2.}
\emph{For $u \in C^2(\Omega) \cap C^0(\overline{\Omega})$ we have}

\begin{equation}
\supU u \leq \supBd u + \frac{d}{\omega_n^{1/n}}\left(\int_{\Gamma^+} \abs{\detm D^2u}\right)^{1/n}
\tag{9.8}
\end{equation}

\emph{where $d = \diam \Omega$.}

\noindent\emph{Proof.} By replacing $u$ with $u - \sup_{\partial\Omega} u$, it suffices to assume $u \leq 0$ on $\partial\Omega$. The $n$-dimensional Lebesgue measure of the normal image of $\Omega$ is given by

\begin{equation}
\abs{\chi(\Omega)} = \abs{\chi(\Gamma^+)}
= \abs{Du(\Gamma^+)}
\leq \int_{\Gamma^+} \abs{\detm D^2u}
\tag{9.9}
\end{equation}

since $D^2u \leq 0$ on $\Gamma^+$. Formula (9.9) can be realised as a consequence of the classical change of variables formula by considering, for positive $\varepsilon$, the mapping $\chi_\varepsilon = \chi - \varepsilon I$, whose Jacobian matrix $D^2u - \varepsilon I$ is then strictly negative in a neighbourhood of $\Gamma^+$, and by subsequently letting $\varepsilon \to 0$. Furthermore, the mapping $\chi_\varepsilon$ is readily shown to be one-to-one on $\Gamma^+$, so that the equality holds in (9.9).


% PDF page 234
\source{gilbarg-trudinger:page-0234}

Let us now show that $u$ can be estimated in terms of $|\chi(\Omega)|$. Suppose that $u$ takes
a positive maximum at a point $y \in \Omega$, and let $k$ be the function whose graph is the
cone $K$ with vertex $(y,u(y))$ and base $\partial\Omega$. Then $\chi_k(\Omega) \subset \chi_u(\Omega)$ since for each sup-
porting hyperplane to $K$, there exists a parallel hyperplane tangent to the graph of $u$.
Now let $\tilde{k}$ be the function whose graph is the cone $\tilde{K}$ with vertex $(y,u(y))$ and base
$B_d(y)$. Clearly, $\chi_k(\Omega) \subset \chi_{\tilde{k}}(\Omega)$; and, consequently,

\[
|\chi_{\tilde{k}}(\Omega)| \le |\chi_u(\Omega)|.
\]

But then using (9.7) and (9.9), we have

\[
\omega_n \left(\frac{u(y)}{d}\right)^n \le \int_{\Gamma^+} |\det D^2u|
\]

and hence

\[
u(y) \le \frac{d}{\omega_n^{1/n}} \left(\int_{\Gamma^+} |\det D^2u|\right)^{1/n}
\]

as required. $\square$

The special case of the estimate (9.4), when $b = 0$, follows from Lemma 9.2 via
the matrix inequality

\begin{equation}\tag{9.10}
\det A \det B \le \left(\frac{\trace AB}{n}\right)^n,\qquad A,B\ \text{symmetric} \ge 0.
\end{equation}

Taking $A = -D^2u$, $B = [a^{ij}]$, we therefore have on $\Gamma^+$

\[
|\det D^2u| = \det(-D^2u)
\le \frac{1}{\mathscr{D}^*}\left(\frac{-a^{ij}D_{ij}u}{n}\right)^n.
\]

We formulate the resulting estimate as follows for later reference.

\medskip

\textbf{Lemma 9.3.}
\textit{For $u \in C^2(\Omega) \cap C^0(\bar{\Omega})$, we have}

\begin{equation}\tag{9.11}
\sup_{\Omega} u \le \sup_{\partial\Omega} u + \frac{d}{n\omega_n^{1/n}}
\left\|\frac{a^{ij}D_{ij}u}{\mathscr{D}^*}\right\|_{L^n(\Gamma^+)}
\end{equation}

The full estimate (9.4) is derived from the following generalization of Lemmas
9.2 and 9.3.


% PDF page 235
\source{gilbarg-trudinger:page-0235}




\noindent 9.1. Maximum Principles for Strong Solutions

\bigskip

\textbf{Lemma 9.4.}
\textit{Let $g$ be a nonnegative, locally integrable function on $\mathbb{R}^n$. Then for any $u \in C^2(\Omega) \cap C^0(\overline{\Omega})$, we have}

\begin{equation}\tag{9.12}
\int_{B_{\Mtilde}(0)} g \leq \int_{\Gamma^+} g(Du)\lvert \Det D^2u\rvert
\leq \int_{\Gamma^+} g(Du)\left(-\frac{a^{ij}D_{ij}u}{n\mathcal{D}^*}\right)^n
\end{equation}

\noindent where

\[
\Mtilde = (\sup_{\Omega} u - \sup_{\partial\Omega} u)/d, \qquad d = \diam \Omega.
\]

\noindent \textit{Proof.} The proof of Lemma 9.4 follows those of Lemmas 9.2 and 9.3, which correspond to the special case $g \equiv 1$. Instead of (9.9), we have the more general formula

\begin{equation}\tag{9.13}
\int_{\chi_u(\Omega)} g \leq \int_{\Gamma^+} g(Du)\lvert \Det D^2u\rvert
\end{equation}

\noindent and since $\chi_k(\Omega) \subset \chi_u(\Omega)$ the estimate (9.12) follows by virtue of (9.7) and (9.10). \hfill $\square$

\medskip

\noindent Let us now suppose that $u \in C^0(\overline{\Omega}) \cap C^2(\Omega)$ and satisfies $Lu \geq f$ in $\Omega$ with condition (9.3). For the weight function $g$ we take

\[
g(p) = \left(\lvert p\rvert^{n/(n-1)} + \mu^{n/(n-1)}\right)^{1-n}
\]

\noindent for some $\mu > 0$, to be fixed later. Then using H\"older's inequality, we have, in
$\Omega^+ = \{x \in \Omega \mid u(x) > 0\}$,

\[
-\frac{a^{ij}D_{ij}u}{n\mathcal{D}^*}
\leq \frac{b^iD_iu - f}{n\mathcal{D}^*}
\leq \frac{\lvert b\rvert\lvert Du\rvert + \lvert f\rvert}{n\mathcal{D}^*}
\leq \frac{\left(\lvert b\rvert^n + \mu^{-n}\lvert f\rvert^n\right)^{1/n}}{n g^{1/n}\mathcal{D}^*}.
\]

\noindent Hence, by (9.12)

\[
\int_{B_{\Mtilde}} g \leq \frac{1}{n^n}\int_{\Gamma^+} \left(\lvert b\rvert^n + \mu^{-n}\lvert f\rvert^n\right)/\mathcal{D}.
\]


% PDF page 236
\source{gilbarg-trudinger:page-0236}

The integral on the left-hand side can be estimated using a further consequence of Hölder's inequality,

\[
g(p) \ge 2^{2-n} \bigl(|p|^n+\mu^n\bigr)^{-1}.
\]

Accordingly, we obtain

\[
\omega_n \log \left(\frac{\wt{M}^n}{\mu^n}+1\right)
\le \frac{2^{n-2}}{n^n}\int_{\Gamma^+}\bigl(|b|^n+\mu^{-n}|f|^n\bigr)\,\D.
\]

If $f \not\equiv 0$, we choose $\mu=\norm{f/\D^*}_{L^n(\Gamma^+)}$ thus to obtain

\[
(9.14)\qquad
\wt{M}\le
\left\{
\exp\left[
\frac{2^{n-2}}{n^n\omega_n}\int_{\Gamma^+}\left(1+\frac{|b|^n}{\D}\right)
-1
\right]
\right\}^{1/n}
\norm{f/\D^*}_{L^n(\Gamma^+)},
\]

while for $f\equiv 0$, we let $\mu\to 0$ so that (9.14) is again satisfied.
This establishes the estimate (9.4) for functions $u\in C^0(\overline{\Omega})\cap C^2(\Omega)$.
The extension to functions $u\in C^0(\overline{\Omega})\cap W^{2,n}_{\operatorname{loc}}(\Omega)$ can be carried out by an approximation argument.
Suppose first that $L$ is uniformly elliptic in $\Omega$ with the ratio $|b|/\lambda$ also bounded.
Let $\{u_m\}$ be a sequence of functions in $C^2(\Omega)$ converging in the sense of $W^{2,n}_{\operatorname{loc}}(\Omega)$ to $u$.
For arbitrary $\varepsilon>0$, we can then assume that $u_m$ converges to $u$ in $W^{2,n}_{\operatorname{loc}}(\Omega_\varepsilon)$ and $u_m\le \varepsilon+\sup_{\partial\Omega_\varepsilon}u$ for some domain $\Omega_\varepsilon\Subset\Omega$.
Consequently, by applying (9.4) to the functions $u_m$ (with $\Omega$ replaced by $\Omega^+$), we obtain

\[
\sup_{\Omega_\varepsilon} u_m
\le \varepsilon+\sup_{\partial\Omega}u^+
+\frac{C}{\lambda}\norm{a^{ij}D_{ij}(u_m-u)+b^iD_i(u_m-u)}_{L^n(\Omega_\varepsilon)}
+C\norm{f/\D^*}_{L^n(\Omega_\varepsilon)},
\]

and hence, letting $m\to\infty$ and using the fact that $\{u_m\}$ converges uniformly to $u$ on $\Omega_\varepsilon$, we have

\[
(9.15)\qquad
\sup_{\Omega_\varepsilon}u
\le \varepsilon+\sup_{\partial\Omega}u^+
+C\norm{f/\D^*}_{L^n(\Omega_\varepsilon)}
\]

from which (9.4) follows as $\varepsilon\to 0$.
To remove the above restrictions on $L$, we consider for $\eta>0$, the operators

\[
L_\eta=\eta(1+|b|)\Delta+L.
\]

We obtain accordingly from (9.15),

\[
\sup_{\Omega_\varepsilon}u
\le \varepsilon+\sup_{\partial\Omega}u^+
+C\left\{
\left\|
\frac{\eta(1+|b|)\Delta u}{\D^*_\eta}
\right\|_{L^n(\Omega_\varepsilon)}
+\left\|\frac{f}{\D^*}\right\|_{L^n(\Omega_\varepsilon)}
\right\}.
\]


% PDF page 237
\source{gilbarg-trudinger:page-0237}

so that letting $\eta \to 0$ and using the dominated convergence theorem, we get
inequality (9.15) again. Theorem 9.1 finally follows by letting $\epsilon \to 0$. \hfill$\square$

\medskip

Note that the estimate (9.14) does not yield (9.11) when $b = 0$. In fact, the
constant in (9.11) can be improved, and a corresponding sharp improvement of (9.14)
can be obtained by explicitly integrating the function $g$ and optimizing the choice of
$\mu$; (see Problems 9.1, 9.2). Also the dependence on diam $\Omega$ can be replaced by a
dependence on $|\hat{\Omega}|$, where $\hat{\Omega}$ is the convex hull of $\Omega$; (see Problem 9.3).
The case $f \equiv 0$ in Theorem 9.1 yields an extension of the weak maximum
principles, Theorem 3.1 and Corollary 3.2. The following uniqueness result for the
Dirichlet problem for strong solutions which extends Theorem 3.3, also follows
automatically.

\medskip

\textbf{Theorem 9.5.}
\emph{Let $L$ be elliptic in the bounded domain $\Omega$ and satisfy (9.3).
Suppose that $u$ and $v$ are functions in $W^{2,n}_{\mathrm{loc}}(\Omega)\cap C^0(\bar{\Omega})$ satisfying
$Lu = Lv$ in $\Omega$, $u = v$ on $\partial\Omega$. Then $u = v$ in $\Omega$.}

\medskip

We can also deduce from the weak maximum principle a generalization of the
strong maximum principle, Theorem 3.5. As in the hypotheses of Theorem 3.5, we
assume that the operator $L$ is uniformly elliptic in $\Omega$ and that $|b|/\lambda, c/\lambda$ are bounded.

\medskip

\textbf{Theorem 9.6.}
\emph{If $u \in W^{2,n}_{\mathrm{loc}}(\Omega)$ satisfies $Lu \geq 0$ in $\Omega$ and
$c = 0$ ($c \leq 0$), then $u$ cannot achieve a maximum (nonnegative maximum) in $\Omega$ unless
it is a constant.}

\medskip

\noindent\textbf{Proof.} If $u$ is differentiable we can follow the proof of Theorem 3.5 with Theorem
9.1 being used in place of Corollary 3.2. In general a slight modification of the proof
of Theorem 3.5 suffices. If we assume, contrary to the theorem, that $u$ is nonconstant
in $\Omega$ and assumes its maximum $M$ in $\Omega$, there must exist concentric balls
$B_\rho(y)\subset B_R(y)\subset \Omega$ such that $u < M$ in $\bar{B}_\rho(y)$ and $u(x_0)=M$ for some $x_0\in B_R(y)$.
But then using Theorem 9.1 and the auxiliary function $v$ defined in the proof of
Lemma 3.4 with $v(x_0)=0$, we have $M-u-\epsilon v>0$, for some $\epsilon>0$, in the annular
region $A=B_R(y)-B_\rho(y)$, which is a contradiction at $x=x_0$. \hfill$\square$

\medskip

\noindent\textbf{9.2. \ $L^p$ Estimates: Preliminary Analysis}

\medskip

Our route to the basic $L^p$ estimates of this chapter is via interpolation. In this sec-
tion, we develop some preliminary analysis: namely, a cube decomposition pro-
cedure, also necessary for the H\"older estimates in Section 9.7; and the Marcin-
kiewicz interpolation theorem that is applied in the next section.

% PDF page 238
\source{gilbarg-trudinger:page-0238}

Cube Decomposition

Let $K_0$ be a cube in $\mathbb{R}^n$, $f$ a nonnegative integrable function defined in $K_0$ and $t$ a positive number satisfying

\[
\int_{K_0} f \leq t|K_0|.
\]

By bisection of the edges of $K_0$, we subdivide $K_0$ into $2^n$ congruent subcubes with disjoint interiors. Those subcubes $K$, which satisfy

\begin{equation}
\int_K f \leq t|K|,
\tag{9.16}
\end{equation}

are similarly subdivided and the process is repeated indefinitely. Let $\mathcal{S}$ denote the set of subcubes $K$ thus obtained that satisfy

\[
\int_K f > t|K|,
\]

and for each $K \in \mathcal{S}$ denote by $\widetilde{K}$ the subcube whose subdivision gives $K$. Since $|\widetilde{K}|/|K| = 2^n$, we have for any $K \in \mathcal{S}$,

\begin{equation}
t < \frac{1}{|K|}\int_K f \leq 2^n t.
\tag{9.17}
\end{equation}

Furthermore, setting $F = \bigcup_{K \in \mathcal{S}} K$ and $G = K_0 - F$, we have

\begin{equation}
f \leq t \quad \text{a.e. in } G.
\tag{9.18}
\end{equation}

The last inequality (9.18) is a consequence of Lebesgue's differentiation theorem [SN], as each point of $G$ lies in a nested sequence of cubes satisfying (9.16) with diameters tending to zero.

For the pointwise estimates in Section 9.7, we also need to consider the set

\[
\widetilde{F} = \bigcup_{K \in \mathcal{S}} \widetilde{K}
\]

which satisfies, by (9.16),

\begin{equation}
\int_{\widetilde{F}} f \leq t|\widetilde{F}|.
\tag{9.19}
\end{equation}

% PDF page 239
\source{gilbarg-trudinger:page-0239}

In particular, when $f$ is the characteristic function $\chi_{\Gamma}$ of a measurable subset $\Gamma$ of
$K_{0}$, we obtain from (9.18) and (9.19) that

\[
(9.20)\qquad |\Gamma| = |\Gamma \cap \widetilde{F}| \leq t|\widetilde{F}|.
\]

\bigskip

\centerline{\bf 9.3. The Marcinkiewicz Interpolation Theorem}

\medskip

Let $f$ be a measurable function on a domain $\Omega$ (bounded or unbounded) in $\mathbb{R}^{n}$. The
\emph{distribution function} $\mu = \mu_{f}$ of $f$ is defined by

\[
(9.21)\qquad \mu(t) = \mu_{f}(t) = |\{x \in \Omega \mid f(x) > t\}|
\]

for $t > 0$, and measures the relative size of $f$. Note that $\mu$ is a decreasing function on
$(0, \infty)$. The basic properties of the distribution function are embodied in the following
lemma.

\medskip

\noindent{\bf Lemma 9.7.} \emph{For any $p > 0$ and $|f|^{p} \in L^{1}(\Omega)$, we have}

\[
(9.22)\qquad \mu(t) \leq t^{-p}\int_{\Omega} |f|^{p},
\]

\[
(9.23)\qquad \int_{\Omega} |f|^{p} = p\int_{0}^{\infty} t^{p-1}\mu(t)\,dt.
\]

\noindent{\it Proof.} Clearly,

\[
\mu(t)t^{p} \leq \int_{|f|\geq t} |f|^{p}
\]

\[
\leq \int_{\Omega} |f|^{p}
\]

for all $t > 0$ and hence (9.22) follows. Next, suppose that $f \in L^{1}(\Omega)$. Then, by
Fubini’s theorem,

\[
\int_{\Omega} |f| = \int_{\Omega}\int_{0}^{|f(x)|} dt\,dx
\]

\[
= \int_{0}^{\infty}\mu(t)\,dt,
\]

and the result (9.23) for general $p$ follows by a change of variables. \hfill$\square$

% PDF page 240
\source{gilbarg-trudinger:page-0240}

We prove the Marcinkiewicz interpolation theorem in the following restricted form:

\textbf{Theorem 9.8.}
\emph{Let $T$ be a linear mapping from $L^q(\Omega)\cap L^r(\Omega)$ into itself, $1 \le q < r < \infty$ and suppose there are constants $T_1$ and $T_2$ such that}

\begin{equation}\tag{9.24}
\mu_{Tf}(t) \le \left(\frac{T_1\|f\|_q}{t}\right)^q,\qquad
\mu_{Tf}(t) \le \left(\frac{T_2\|f\|_r}{t}\right)^r
\end{equation}

\emph{for all $f \in L^q(\Omega)\cap L^r(\Omega)$ and $t > 0$. Then $T$ extends as a bounded linear mapping from $L^p(\Omega)$ into itself for any $p$ such that $q < p < r$, and}

\begin{equation}\tag{9.25}
\|Tf\|_p \le C T_1^{\alpha} T_2^{1-\alpha} \|f\|_p
\end{equation}

\emph{for all $f \in L^q(\Omega)\cap L^r(\Omega)$, where}

\[
\frac{1}{p} = \frac{\alpha}{q} + \frac{1-\alpha}{r}
\]

\emph{and $C$ depends only on $p$, $q$ and $r$.}

\textbf{Proof.} \emph{For $f \in L^q(\Omega)\cap L^r(\Omega)$ and $s > 0$, we write}

\[
f = f_1 + f_2,
\]

\emph{where}

\[
f_1(x)=\begin{cases}
f(x) & \text{if } |f(x)| > s,\\
0 & \text{if } |f(x)| \le s,
\end{cases}
\qquad
f_2(x)=\begin{cases}
0 & \text{if } |f(x)| > s,\\
f(x) & \text{if } |f(x)| \le s.
\end{cases}
\]

\emph{Then}

\[
|Tf| \le |Tf_1| + |Tf_2|,
\]

\emph{and hence}

\[
\mu(t)=\mu_{Tf}(t) \le \mu_{Tf_1}(t/2)+\mu_{Tf_2}(t/2).
\]

\[
\le \left(\frac{2T_1}{t}\right)^q \int_{\Omega} |f_1|^q + \left(\frac{2T_2}{t}\right)^r \int_{\Omega} |f_2|^r.
\]


% PDF page 241
\source{gilbarg-trudinger:page-0241}

Therefore, by Lemma 9.7 we have

\[
\int_{\Omega} |Tf|^p = p \int_0^{\infty} t^{p-1}\mu(t)\,dt
\]

\[
\leq p(2T_1)^q \int_0^{\infty} t^{p-1-q}
\left(\int_{|f|>s} |f|^q\right)\,dt
\]

\[
\quad + p(2T_2)^r \int_0^{\infty} t^{p-1-r}
\left(\int_{|f|\leq s} |f|^r\right)\,dt
\]

Let us now choose $s$ as a function of $t$; in particular, we take $t = As$ for some positive number $A$ to be fixed later. Then, we have

\[
\int_{\Omega} |Tf|^p \leq p(2T_1)^q A^{p-q} \int_0^{\infty} s^{p-1-q}
\left(\int_{|f|>s} |f|^q\right)\,ds
\]

\[
\quad + p(2T_2)^r A^{p-r} \int_0^{\infty} s^{p-1-r}
\left(\int_{|f|\leq s} |f|^r\right)\,ds.
\]

But

\[
\int_0^{\infty} s^{p-1-q}\left(\int_{|f|>s} |f|^q\right)\,ds
= \int_{\Omega} |f|^q \left(\int_0^{|f|} s^{p-1-q}\,ds\right)
\]

\[
= \frac{1}{p-q}\int_{\Omega} |f|^p,
\]

and, similarly,

\[
\int_0^{\infty} s^{p-1-r}\left(\int_{|f|\leq s} |f|^r\right)\,ds
= \int_{\Omega} |f|^r \left(\int_{|f|}^{\infty} s^{p-1-r}\,ds\right)
\]

\[
= \frac{1}{r-p}\int_{\Omega} |f|^p.
\]

% PDF page 242
\source{gilbarg-trudinger:page-0242}
\section*{9.4. The Calderon-Zygmund Inequality}

In this section we establish the basic $L^p$ estimates for Poisson's equation through a
further consideration of the Newtonian potential, previously treated in Chapter 4.
Let $\Omega$ be a bounded domain in $\mathbb{R}^n$ and $f$ a function in $L^p(\Omega)$ for some $p \ge 1$. Recall
that the Newtonian potential of $f$ is the function $w = Nf$ defined by the convolution.

\begin{equation}
\tag{9.26}
w(x)=\int_{\Omega}\Gamma(x-y)f(y)\,dy,
\end{equation}

where $\Gamma$ is the fundamental solution of Laplace's equation given by (4.1). The
following result, which embraces a special case of the Calderon-Zygmund inequality, is the $L^p$ analogue of the Hölder estimate of Lemma 4.4.

\textbf{Theorem 9.9.}
\emph{Let $f \in L^p(\Omega)$, $1 < p < \infty$, and let $w$ be the Newtonian potential of $f$.
Then $w \in W^{2,p}(\Omega)$, $\Delta w = f$ a.e. and}
\begin{equation}
\tag{9.27}
\|D^2w\|_p \le C\|f\|_p
\end{equation}

\emph{where $C$ depends only on $n$ and $p$. Furthermore, when $p = 2$ we have}

\begin{equation}
\tag{9.28}
\int_{\mathbb{R}^n}\!|D^2w|^2=\int_{\Omega}\!f^2
\end{equation}

% PDF page 243
\source{gilbarg-trudinger:page-0243}

\pagestyle{empty}


\section*{9.4. The Calderon-Zygmund Inequality}

\textit{Proof.} \textbf{(i)} Let us deal first with the case $p = 2$. If $f \in C^\infty_0(\mathbb{R}^n)$, we have $w \in C^\infty(\mathbb{R}^n)$ and, by Lemma 4.3, $\Delta w = f$. Consequently, for any ball $B_R$ containing the support of $f$,

\[
\int_{B_R} (\Delta w)^2 = \int_{B_R} f^2.
\]

Applying Green's first identity (2.10) twice, we then obtain

\[
\int_{B_R} |D^2 w|^2 = \int_{B_R} \sum (D_{ij} w)^2
\]

\[
= \int_{B_R} f^2 + \int_{\partial B_R} Dw \cdot \frac{\partial}{\partial v} Dw.
\]

Using (2.14) we have

\[
Dw = O(R^{1-n}),\ D^2 w = O(R^{-n})
\]

uniformly on $\partial B_R$ as $R \to \infty$, whence the identity (9.28) follows. To extend (9.28) to arbitrary $f \in L^2(\Omega)$, we observe first that, by Lemma 7.12, $N$ is a bounded mapping from $L^p$ into itself for $1 \leq p < \infty$. The full strength of Theorem 9.9 in the case $p = 2$ then follows by approximation. Indeed, if the sequence $\{f_m\} \subset C^\infty_0(\Omega)$ converges to $f$ in $L^2(\Omega)$, the sequence of Newtonian potentials $\{Nf_m\}$ converges to $w$ in $W^{2,2}(\Omega)$.

\textbf{(ii)} For fixed $i, j$, we now define the linear operator $T: L^2(\Omega) \to L^2(\Omega)$ by

\[
Tf = D_{ij} w.
\]

From Lemma 9.7 and (9.28), we have

\begin{equation}
\mu(t) = \mu_{Tf}(t) \leq \left(\frac{\|f\|_2}{t}\right)^2
\tag{9.29}
\end{equation}

for all $t > 0$ and $f \in L^2(\Omega)$. We now show that, in addition,

\begin{equation}
\mu(t) \leq C \frac{\|f\|_1}{t}
\tag{9.30}
\end{equation}


% PDF page 244
\source{gilbarg-trudinger:page-0244}


for all $t>0$ and $f\in L^2(\Omega)$, thereby making possible the application of the Marcinkiewicz interpolation theorem. To accomplish this we first extend $f$ to vanish outside $\Omega$ and fix a cube $K_0 \supset \Omega$, so that for fixed $t>0$ we have

\[
\int_{K_0} f \le t|K_0|.
\]

The cube $K_0$ is now decomposed according to the procedure described in Section 9.2 yielding a sequence of parallel subcubes $\{K_l\}_{l=1}^{\infty}$ such that

\begin{equation}
t < \frac{1}{|K_l|}\int_{K_l} |f| < 2^n t
\tag{9.31}
\end{equation}

and

\[
|f|\le t \qquad \text{a.e. on } G = K_0 - \bigcup K_l.
\]

The function $f$ is now split into a ``good part'' $g$ defined by

\[
g(x)=
\begin{cases}
f(x) & \text{for } x\in G \\
\frac{1}{|K_l|}\int_{K_l} f & \text{for } x\in K_l,\quad l=1,2,\ldots
\end{cases}
\]

and a ``bad part'' $b=f-g$. Clearly,

\[
|g|\le 2^n t \qquad \text{a.e.,}
\]

\[
b(x)=0 \qquad \text{for } x\in G,
\]

\[
\int_{K_l} b=0 \qquad \text{for } l=1,2,\ldots .
\]

Since $T$ is linear, $Tf=Tg+Tb$, and, hence,

\[
\muTf(t)\le \muTg(t/2)+\muTb(t/2).
\]

\noindent (iii) \textit{Estimation of $Tg$:} By (9.29)

\[
\muTg(t/2)\le \frac{4}{t^2}\int g^2
\]

\[
\le \frac{2^{n+2}}{t}\int |g|
\]

\[
\le \frac{2^{n+2}}{t}\int |f|
\]


% PDF page 245
\source{gilbarg-trudinger:page-0245}

9.4. The Calderon-Zygmund Inequality \hfill 233

(iv) \textit{Estimation of $Tb$: Writing}

\[
b_l = b\chi_{K_l} =
\begin{cases}
b & \text{on } K_l,\\
0 & \text{elsewhere,}
\end{cases}
\]

we have

\[
Tb = \sum_{l=1}^{\infty} Tb_l.
\]

Let us now fix some $l$ and a sequence $\{b_{lm}\} \subset C_0^{\infty}(K_l)$ converging to $b_l$ in $L^2(\Omega)$ and
satisfying

\[
\int_{K_l} b_{lm} = \int_{K_l} b_l
\]

\[
= 0.
\]

Then for $x \notin K_l$, we have the formula

\[
Tb_{lm}(x) = \int_{K_l} D_{ij}\Gamma(x-y)b_{lm}(y)\,dy
\]

\[
= \int_{K_l} \{D_{ij}\Gamma(x-y)-D_{ij}\Gamma(x-\bar y)\}b_{lm}(y)\,dy,
\]

where $\bar y = \bar y_l$ denotes the center of $K_l$. Letting $\delta = \delta_l$ denote the diameter of $K_l$, we
then obtain (with an estimation similar to that of the integral $I_6$ in the proof of
Lemma 4.4),

\[
|Tb_{lm}(x)| \le C(n)\delta[\mathrm{dist}(x,K_l)]^{-n-1}\int_{K_l}|b_{lm}(y)|\,dy.
\]

Letting $B_l = B_{\delta}(\bar y)$ denote the concentric ball of radius $\delta$, we obtain by integration

\[
\int_{K_0 - B_l} |Tb_{lm}| \le C(n)\delta \int_{|x|\ge \delta/2}\frac{dx}{|x|^{n+1}}\int_{K_l}|b_{lm}|
\]

\[
\le C(n)\int_{K_l}|b_{lm}|.
\]


% PDF page 246
\source{gilbarg-trudinger:page-0246}

Consequently, letting $m \to \infty$, writing $F^* = \bigcup B_l$, $G^* = K_0 - F^*$ and summing
over $l$, we get

\[
\int_{G^*} |Tb| \le C(n) \int |b|
\]

\[
\le C(n) \int |f|,
\]

so that by Lemma 9.7

\[
\left|\{x \in G^* \mid |Tb| > t/2\}\right| \le \frac{C\|f\|_1}{t}.
\]

However, by (9.31),

\[
|F^*| \le \omega_n n^{n/2}|F|
\]

\[
\le \frac{C\|f\|_1}{t},
\]

and thus (9.30) holds.

\noindent (v) To conclude the proof of Theorem 9.9, we note that (9.29) and (9.30)
fulfill the hypotheses of the Marcinkiewicz interpolation theorem (Theorem 9.8)
with $q = 1$, $r = 2$. Consequently,

\begin{equation}
\|Tf\|_p \le C(n, p)\|f\|_p
\tag{9.32}
\end{equation}

for all $1 < p \le 2$ and $f \in L^2(\Omega)$. The inequality (9.32) is extended to $p > 2$ by
duality. For if $f, g \in C_0^{\infty}(\Omega)$, then

\[
\int_{\Omega} (Tf)g = \int_{\Omega} wD_{ij}g
\]

\[
= \int\!\!\int_{\Omega\Omega} \Gamma(x-y)f(y)D_{ij}g(x)\,dx\,dy
\]

\[
= \int_{\Omega} f\,Tg
\]

\[
\le \|f\|_p\|Tg\|_{p'}.
\]


% PDF page 247
\source{gilbarg-trudinger:page-0247}

\begin{quote}
Thus if $p>2$, we have from (9.32),
\[
\|Tf\|_p=\sup\left\{\int_\Omega (Tf)g \,\|g\|_{p'}=1\right\}
\]
\[
\le C(n,p')\|f\|_p,
\]
so that (9.32) holds for all $1<p<\infty$. As in the case $p=2$, we can then infer the full conclusion of Theorem 9.9 by approximation. $\square$

Note that $T$ can be defined as a bounded operator on $L^p(\Omega)$ even if $\Omega$ is unbounded, in which case the conclusion of Theorem 9.9 still holds provided $n\ge 3$. Other approaches to the derivation of inequality (9.27) are discussed in the Notes to this chapter.

The $L^p$ estimates for solutions of Poisson's equation follow immediately from Theorem 9.9.

\textbf{Corollary 9.10.}
\emph{Let $\Omega$ be a domain in $\mathbb{R}^n$, $u\in W^{2,p}_0(\Omega)$, $1<p<\infty$. Then}
\begin{equation}\tag{9.33}
\|D^2u\|_p\le C\|\Delta u\|_p
\end{equation}
\emph{where $C=C(n,p)$. If $p=2$,}
\begin{equation}\tag{9.34}
\|D^2u\|_2=\|\Delta u\|_2.
\end{equation}
\end{quote}

\section*{9.5. \emph{$L^p$ Estimates}}

In this section, we derive interior and global $L^p$ estimates for the second derivatives of elliptic equations of the forms (9.1) and (9.2). The technique of perturbation from the constant coefficient case is similar to that used for the derivation of the Schauder estimates in Sections 6.1 and 6.2. We first deal with interior estimates; the following theorem is analogous to Theorem 6.1.

\textbf{Theorem 9.11.}
\emph{Let $\Omega$ be an open set in $\mathbb{R}^n$ and $u\in W^{2,p}_{\mathrm{loc}}(\Omega)\cap L^p(\Omega)$, $1<p<\infty$, a strong solution of the equation $Lu=f$ in $\Omega$ where the coefficients of $L$ satisfy, for positive constants $\lambda,\Lambda$,}
\[
a^{ij}\in C^0(\Omega), \qquad b^i,c\in L^\infty(\Omega), \qquad f\in L^p(\Omega);
\]
\begin{equation}\tag{9.35}
a^{ij}\xi_i\xi_j\ge \lambda |\xi|^2 \qquad \forall \xi\in \mathbb{R}^n;
\end{equation}
\[
|a^{ij}|, |b^i|, |c| \le \Lambda,
\]
\emph{where $i,j=1,\ldots,n$. Then for any domain $\Omega' \Subset \Omega$,}
\begin{equation}\tag{9.36}
\|u\|_{2,p;\Omega'}\le C(\|u\|_{p;\Omega}+\|f\|_{p;\Omega}),
\end{equation}

% PDF page 248
\source{gilbarg-trudinger:page-0248}

\[
\text{where } C \text{ depends on } n, p, \lambda, \Lambda, \Omega', \Omega \text{ and the moduli of continuity of the coefficients}
\]
\[
a^{ij} \text{ on } \Omega'.
\]

\textit{Proof.} For a fixed point $x_0 \in \Omega'$, we let $L_0$ denote the constant coefficient operator
given by

\[
L_0 u = a^{ij}(x_0)D_{ij}u.
\]

By means of the linear transformation $Q$ used in the proof of Lemma 6.1, we obtain
from Corollary 9.10, the estimate

\[
\tag{9.37}
\|D^2 v\|_{p;\Omega} \le \frac{C}{\lambda}\|L_0 v\|_{p;\Omega}
\]

for any $v \in W^{2,p}_0(\Omega)$, where $C = C(n,p)$ as in (9.33). Consequently, if $v$ has support
in a ball $B_R = B_R(x_0) \subset \Omega$, we have

\[
L_0v = \bigl(a^{ij}(x_0)-a^{ij}\bigr)D_{ij}v + a^{ij}D_{ij}v,
\]

and by (9.37)

\[
\|D^2 v\|_p \le \frac{C}{\lambda}\left(\sup_{B_R}\lvert a-a(x_0)\rvert \,\|D^2 v\|_p + \|a^{ij}D_{ij}v\|_p\right),
\]

where $a = [a^{ij}]$. Since $a$ is uniformly continuous on $\Omega'$, there exists a positive
number $\delta$ such that

\[
\lvert a-a(x_0)\rvert \le \lambda/2C
\]

if $\lvert x-x_0\rvert < \delta$, and hence

\[
\|D^2 v\|_p \le C\|a^{ij}D_{ij}v\|_p
\]

provided $R \le \delta$, where $C = C(n,p,\lambda)$.

For $\sigma \in (0,1)$, we now introduce a cutoff function $\eta \in C^2(B_R)$ satisfying $0 \le
\eta \le 1$, $\eta = 1$ in $B_{\sigma R}$, $\eta = 0$ for $\lvert x\rvert \ge \sigma'R$, $\sigma' = (1+\sigma)/2$, $\lvert D\eta\rvert \le 4/((1-\sigma)R)$,
$\lvert D^2\eta\rvert \le 16/((1-\sigma)^2R^2)$. Then, if $u \in W^{2,p}_{\mathrm{loc}}(\Omega)$ satisfies $Lu = f$ in $\Omega$ and $v = \eta u$, we
obtain

\[
\|D^2 u\|_{p;B_{\sigma R}} \le C\|\eta a^{ij}D_{ij}u + 2a^{ij}D_i\eta D_j u + ua^{ij}D_{ij}\eta\|_{p;B_R}
\]
\[
\le C\left(\|f\|_{p;B_R} + \frac{1}{(1-\sigma)R}\|Du\|_{p;B_{\sigma'R}}
+ \frac{1}{(1-\sigma)^2R^2}\|u\|_{p;B_R}\right)
\]

provided $R \le \delta \le 1$, where $C = C(n,p,\lambda,\Lambda)$.

% PDF page 249
\source{gilbarg-trudinger:page-0249}

Introducing the weighted seminorms

\[
\Phi_k = \sup_{0<\sigma<1}(1-\sigma)^kR^k\|D^ku\|_{p;B_{\sigma R}}, \qquad k=0,1,2,
\]

we, therefore, have

\[
\tag{9.38}
\Phi_2 \leq C(R^2\|f\|_{p;B_R}+\Phi_1+\Phi_0).
\]

We claim now that $\Phi_k$ satisfy an interpolation inequality

\[
\tag{9.39}
\Phi_1 \leq \varepsilon \Phi_2 + \frac{C}{\varepsilon}\Phi_0
\]

for any $\varepsilon > 0$, where $C=C(n)$. By its invariance under coordinate stretching it
suffices to prove (9.39) for the case $R=1$.

For $\gamma > 0$, we fix $\sigma=\sigma_\gamma$ so that

\[
\Phi_1 \leq (1-\sigma_\gamma)\|Du\|_{p;B_\sigma}+\gamma
\]
\[
\leq \varepsilon(1-\sigma)^2\|D^2u\|_{p;B_\sigma}+\frac{C}{\varepsilon}\|u\|_{p;B_\sigma}+\gamma
\]

by Theorem 7.28, so that letting $\gamma\to 0$, we obtain (9.39). Using (9.39) in (9.38), we
then get

\[
\Phi_2 \leq C(R^2\|f\|_{p;B_R}+\Phi_0),
\]

that is,

\[
\tag{9.40}
\|D^2u\|_{p;B_{\sigma R}} \leq \frac{C}{(1-\sigma)^2R^2}\left(R^2\|f\|_{p;B_R}+\|u\|_{p;B_R}\right),
\]

where $C=C(n,p,\lambda,\Lambda)$ and $0<\sigma<1$.
The desired estimate (9.36) follows by taking $\sigma=\frac{1}{2}$ and covering $\Omega'$ with a
finite number of balls of radius $R/2$ for $R\leq \min\{\delta,\operatorname{dist}(\Omega',\partial\Omega)\}$. \(\square\)

In order to extend Theorem 9.11 to the boundary $\partial\Omega$, we first consider the case
of a flat boundary portion. Letting

\[
\Omega^+ = \Omega \cap \mathbb{R}_+^n = \{x\in \Omega\mid x_n > 0\},
\]
\[
(\partial\Omega)^+ = (\partial\Omega)\cap \mathbb{R}_+^n = \{x\in \partial\Omega\mid x_n > 0\},
\]

we have the following extension of Corollary 9.10.

% PDF page 250
\source{gilbarg-trudinger:page-0250}

\noindent 238 \hfill 9. Strong Solutions

\noindent \textbf{Lemma 9.12.} \textit{Let $u \in W^{1,1}_{0}(\Omega^+)$, $f \in L^p(\Omega^+)$, $1 < p < \infty$, satisfy $\Delta u = f$ weakly in $\Omega^+$ with $u = 0$ near $(\partial\Omega)^+$. Then $u \in W^{2,p}(\Omega^+) \cap W^{1,p}_0(\Omega^+)$ and}

\begin{equation}
\tag{9.41}
\|D^2u\|_{p;\Omega^+} \le C\|f\|_{p;\Omega^+},
\end{equation}

\noindent where $C = C(n,p)$.

\noindent \textit{Proof.} We extend $u$ and $f$ to all of $\mathbb{R}^n$ by setting $u = f = 0$ in $\mathbb{R}^n_+ - \Omega$, and then to all of $\mathbb{R}^n$ by odd reflection, that is, by setting

\begin{equation*}
u(x',x_n) = -u(x',-x_n), \qquad f(x',x_n) = -f(x',-x_n)
\end{equation*}

\noindent for $x_n < 0$, where $x' = (x_1,\ldots,x_{n-1})$. It follows that the extended functions satisfy $\Delta u = f$ weakly in $\mathbb{R}^n$. To show this we take an arbitrary test function $\varphi \in C^1_0(\mathbb{R}^n)$, and for $\varepsilon > 0$ let $\eta$ be an even function in $C^1(\mathbb{R})$ such that $\eta(t) = 0$ for $|t| \le \varepsilon$, $\eta(t) = 1$ for $|t| \ge 2\varepsilon$ and $|\eta'| \le 2/\varepsilon$. Then

\begin{equation*}
- \int \eta f\varphi = \int Du \cdot D(\eta\varphi)
\end{equation*}
\begin{equation*}
= \int \eta Du \cdot D\varphi + \int \varphi\eta' D_n u.
\end{equation*}

\noindent Now

\begin{equation*}
\left|\int \varphi\eta' D_n u\right| =
\left|\int_{0 < x_n < 2\varepsilon} \bigl(\varphi(x',x_n) - \varphi(x',-x_n)\bigr)\eta' D_n u\right|
\end{equation*}
\begin{equation*}
\le 8 \max |D\varphi| \int_{0 < x_n < 2\varepsilon} |D_n u|
\end{equation*}
\begin{equation*}
\to 0 \qquad \text{as } \varepsilon \to 0.
\end{equation*}

\noindent Consequently, letting $\varepsilon \to 0$, we obtain

\begin{equation*}
- \int f\varphi = \int Du \cdot D\varphi,
\end{equation*}

\noindent so that $u \in W^{1,1}(\mathbb{R}^n)$ is a weak solution of $\Delta u = f$.
Since $u$ also has compact support in $\mathbb{R}^n$, the regularization $u_h \in C^\infty_0(\mathbb{R}^n)$, and satisfies $\Delta u_h = f_h$ in $\mathbb{R}^n$. Hence, by Lemma 7.2 and Corollary 9.10, $u_h \to u$ in $W^{2,p}(\mathbb{R}^n)$ as $h \to 0$ and, moreover, $u$ satisfies the estimate (9.33). However, then the estimate (9.41) follows with constant $C$ twice that in (9.33). Since $u_h(x',0) = 0$, we also obtain $u \in W^{1,p}_0(\Omega^+)$. $\square$

\noindent For the global estimate, we require that boundary values are taken on in the sense of $W^{1,p}(\Omega)$. If $T$ is a subset of $\partial\Omega$ and $u \in W^{1,p}(\Omega)$, we say that $u = 0$ on $T$ in the


% PDF page 251
\source{gilbarg-trudinger:page-0251}

sense of $W^{1,p}(\Omega)$ if $u$ is the limit in $W^{1,p}(\Omega)$ of a sequence of functions in $C^1(\Omega)$ vanishing near $T$. For the case $p=2$, this coincides with the definition used in Section 8.10 and, when $u$ is continuous on $T$, is implied by $u$ vanishing on $T$ in the usual pointwise sense. With the aid of Lemma 9.12 we now derive a local boundary estimate.

\textbf{Theorem 9.13.} \textit{Let $\Omega$ be a domain in $\mathbb{R}^n$ with a $C^{1,1}$ boundary portion $T \subset \partial \Omega$. Let $u \in W^{2,p}(\Omega)$, $1<p<\infty$, be a strong solution of $Lu=f$ in $\Omega$ with $u=0$ on $T$, in the sense of $W^{1,p}(\Omega)$, where $L$ satisfies (9.35) with $a^{ij} \in C^0(\Omega \cup T)$. Then, for any domain $\Omega' \Subset \Omega \cup T$,}

\begin{equation}
\tag{9.42}
\|u\|_{2,p;\Omega'} \le C(\|u\|_{p;\Omega} + \|f\|_{p;\Omega})
\end{equation}

\textit{where $C$ depends on $n$, $p$, $\lambda$, $\Lambda$, $T$, $\Omega'$, $\Omega$ and the moduli of continuity of the coefficients $a^{ij}$ on $\Omega'$.}

\textit{Proof.} Since $T \in C^{1,1}$, for each point $x_0 \in T$ there is a neighbourhood $\mathcal{N}=\mathcal{N}_{x_0}$ and a diffeomorphism $\psi=\psi_{x_0}$ from $\mathcal{N}$ onto the unit ball $B=B_1(0)$ in $\mathbb{R}^n$ such that $\psi(\mathcal{N}\cap\Omega) \subset \mathbb{R}^+$, $\psi(\mathcal{N}\cap\partial\Omega) \subset \partial\mathbb{R}^+$, $\psi \in C^{1,1}(\mathcal{N})$, $\psi^{-1} \in C^{1,1}(B)$. As in Lemma 6.5, writing $y=\psi(x)=(\psi_1(x),\ldots,\psi_n(x))$, $\tilde{u}(y)=u(x)$, $x\in \mathcal{N}$, $y\in B$, we have

\[
\tilde{L}\tilde{u}=\tilde{a}^{ij}D_{ij}\tilde{u}+\tilde{b}^iD_i\tilde{u}+\tilde{c}\tilde{u}=\tilde{f}
\]

in $B^+$, where

\[
\tilde{a}^{ij}(y)=\frac{\partial\psi_i}{\partial x_r}\frac{\partial\psi_j}{\partial x_s}a^{rs}(x),\qquad
\tilde{b}^i(y)=\frac{\partial^2\psi_i}{\partial x_r\partial x_s}a^{rs}(x)+\frac{\partial\psi_i}{\partial x_r}b^r(x),
\]
\[
\tilde{c}(y)=c(x),\qquad \tilde{f}(y)=f(x)
\]

so that $\tilde{L}$ satisfies conditions similar to (9.35) with constants $\tilde{\lambda}$, $\tilde{\Lambda}$ depending on $\lambda$, $\Lambda$ and $\psi$. Furthermore, $\tilde{u} \in W^{2,p}(B^+)$, and $\tilde{u}=0$ on $B\cap \partial\mathbb{R}^n_+$ in the sense of $W^{1,p}(B^+)$. We now proceed as in the proof of Theorem 9.11 with the ball $B_R(x_0)$ replaced by the half ball $B_R^+(0) \subset B$ and with Lemma 9.12 used in place of Corollary 9.10. We obtain thus, instead of (9.40), the estimate

\[
\|D^2\tilde{u}\|_{p;B_{\sigma R}^+} \le \frac{C}{(1-\sigma)^2R^2}\{R^2\|\tilde{f}\|_{p;B_R^+}+\|\tilde{u}\|_{p;B_R^+}\}
\]

provided $R \le \delta \le 1$, where $C$ depends on $n$, $p$, $\lambda$, $\Lambda$ and $\psi$; and $\delta$ depends on the moduli of continuity of $a^{ij}$ at $x_0$ and also on $\psi$. Taking $\sigma=\frac12$ and $\mathcal{N}=\mathcal{N}_{x_0}=\psi^{-1}(B_{\delta/2})$ we therefore have, on returning to our original coordinates,

\[
\|D^2u\|_{p;\mathcal{N}} \le C(\|u\|_{p;\mathcal{N}}+\|f\|_{p;\mathcal{N}})
\]

where $C=C(n,p,\lambda,\Lambda,\delta,\psi)$. Finally, by covering $\Omega' \cap T$ with a finite number of such neighbourhoods $\mathcal{N}$ and using also the interior estimate (9.36), we obtain the desired estimate (9.42). $\square$


% PDF page 252
\source{gilbarg-trudinger:page-0252}

\noindent When $T=\partial \Omega$ in Theorem 9.13 we may take $\Omega'=\Omega$ to obtain a global $W^{2,p}(\Omega)$ estimate. This estimate can, in fact, be refined as follows.

\textbf{Theorem 9.14.} \textit{Let $\Omega$ be a $C^{1,1}$ domain in $\mathbb{R}^n$ and suppose the operator $L$ satisfies the conditions (9.35) with $a^{ij}\in C^{0}(\overline{\Omega})$, $i, j = 1, \ldots, n$. Then if $u \in W^{2,p}(\Omega) \cap W^{1,p}_0(\Omega)$, $1 < p < \infty$, we have}

\begin{equation}
\tag{9.43}
\|u\|_{2,p;\Omega} \le C\|Lu - \sigma u\|_{p;\Omega}
\end{equation}

\noindent \textit{for all $\sigma \ge \sigma_0$, where $C$ and $\sigma_0$ are positive constants depending only on $n$, $p$, $\lambda$, $\Lambda$, $\Omega$ and the moduli of continuity of the coefficients $a^{ij}$.}

\noindent \textit{Proof.} We define a domain $\Omega_0$ in $\mathbb{R}^{n+1}(x, t)$ by

\[
\Omega_0 = \Omega \times (-1,1)
\]

\noindent together with the operator $L_0$, given by

\[
L_0 v = Lv + D_{tt}v,
\]

\noindent for $v \in W^{2,p}(\Omega_0)$. Then, if $u \in W^{2,p}(\Omega) \cap W^{1,p}_0(\Omega)$, the function $v$, given by

\[
v(x,t)=u(x)\cos \sigma^{1/2}t,
\]

\noindent belongs to $W^{2,p}(\Omega_0)$ and vanishes on $\partial \Omega \times (-1,1)$ in the sense of $W^{1,p}(\Omega_0)$. Furthermore,

\[
L_0 v = \cos \sigma^{1/2}t(Lu-\sigma u),
\]

\noindent so that by Theorem 9.13 with $\Omega'=\Omega \times (-\varepsilon,\varepsilon)$, $0<\varepsilon\le \frac{1}{2}$, we get

\[
\|D_{tt}v\|_{p;\Omega'} \le C(\|Lu-\sigma u\|_{p;\Omega} + \|u\|_{p;\Omega}),
\]

\noindent where $C$ depends on the quantities listed in the statement of the Theorem. But now, taking $\varepsilon=\pi/3\sigma^{1/2}$, we have

\[
\|D_{tt}v\|_{p;\Omega'} = \sigma\|v\|_{p;\Omega'}
\]
\[
\ge \sigma \cos(\sigma^{1/2}\varepsilon)(2\varepsilon)^{1/p}\|u\|_{p;\Omega}
\]
\[
\ge \frac{1}{2}\left(\frac{2\pi}{3}\right)^{1/p}\sigma^{1-1/2p}\|u\|_{p;\Omega}
\]

\noindent so that if $\sigma$ is sufficiently large

\begin{equation}
\tag{9.44}
\|u\|_{p;\Omega} \le C\|Lu - \sigma u\|_{p;\Omega},
\end{equation}

\noindent and hence (9.43) follows from Theorem 9.13. $\square$


% PDF page 253
\source{gilbarg-trudinger:page-0253}

Note that when $p \ge n$, Theorem 9.14 follows directly from Theorems 9.1 and 9.13. In the Hilbert space approach in Chapter 8, the analogue of Theorem 9.14 is Lemma 8.4 and, indeed, it is possible to deduce the case $p=2$ from Lemma 8.4. We remark also that by invoking the Sobolev imbedding theorem (see, in particular, Corollary 7.11) in the proofs of Theorems 9.12, 9.13 and 9.14, we may weaken the conditions on the lower order coefficients of $L$ to $b^i \in L^q(\Omega)$, $c \in L^r(\Omega)$, where $q > n$ if $p \le n$, $q = p$ if $p > n$, $r > n/2$ if $p \le n/2$, $r = p$ if $p > n/2$.

\section*{9.6. The Dirichlet Problem}

The main objective of this section is the following existence and uniqueness theorem for the Dirichlet problem for strong solutions.

\textbf{Theorem 9.15.} \textit{Let $\Omega$ be a $C^{1,1}$ domain in $\mathbb{R}^n$, and let the operator $L$ be strictly elliptic in $\Omega$ with coefficients $a^{ij} \in C^0(\overline{\Omega})$, $b^i$, $c \in L^\infty$, with $i, j = 1, \ldots, n$ and $c \le 0$. Then, if $f \in L^p(\Omega)$ and $\varphi \in W^{2,p}(\Omega)$, with $1 < p < \infty$, the Dirichlet problem $Lu = f$ in $\Omega$, $u - \varphi \in W_0^{1,p}(\Omega)$ has a unique solution $u \in W^{2,p}(\Omega)$.}

\textit{Proof.} There are various methods for deducing Theorem 9.15 from previous existence theorems in Chapters 4, 6 and 8. For example, the case $p \ge n$ may be derived from either Theorems 6.14 or 8.14 by appropriate approximation (Problem 9.7), or alternatively from the special case of Poisson's equation by the method of continuity (Problem 9.8). Our treatment is based on the case $p=2$ already covered by Theorems 8.9 and 8.12 under stronger coefficient hypotheses. We shall need the following regularity result which in fact is a refinement of Theorems 9.11 and 9.13.

\textbf{Lemma 9.16.} \textit{In addition to the hypotheses of Theorem 9.13, suppose that $f \in L^q(\Omega)$ for some $q \in (p, \infty)$. Then, $u \in W^{2,q}_{\mathrm{loc}}(\Omega \cup T)$, $u = 0$ on $T$ in the sense of $W^{1,q}(\Omega)$, and consequently, $u$ satisfies the estimate (9.42) with $p$ replaced by $q$.}

\textit{Proof.} We first treat the interior case when $T$ is empty. Returning to the proof of Theorem 9.11, we fix a ball $B_R = B_R(x_0)$ and a cutoff function $\eta$, and set $v = \eta u$, $g = a^{ij}D_{ij}v$, so that

\[
L_0v = \bigl(a^{ij}(x_0) - a^{ij}(x)\bigr)D_{ij}v + g.
\]

Since $Lu = f$, it follows from the Sobolev imbedding theorem that $g \in L^r(\Omega)$ where
\[
\frac{1}{r} = \max\left\{ \frac{1}{q}, \left(\frac{1}{p} - \frac{1}{n}\right) \right\}.
\]
By means of the linear transformation $Q$, we can diagonalize the matrix $[a^{ij}(x_0)]$ so that the operator $L_0$ becomes the Laplacian, and hence

\[
\Delta \tilde{v} = \bigl(\delta^{ij} - \tilde{a}^{ij}(x)\bigr)D_{ij}\tilde{v} + \tilde{g},
\]

where $\tilde{v}$, $\tilde{a}^{ij}$, $\tilde{g}$ correspond to $v$, $a^{ij}$, $g$, respectively. By taking the Newtonian potential, we then obtain the equation

\[
\tilde{v} = N\bigl[(\delta^{ij} - \tilde{a}^{ij}(x))D_{ij}\tilde{v}\bigr] + N\tilde{g}.
\]


% PDF page 254
\source{gilbarg-trudinger:page-0254}

Consequently, the function $v$ satisfies an equation of the form

\[
(9.45)\qquad v = Tv + h,
\]

where by virtue of the Calderon-Zygmund estimate (Theorem 9.9) $T$ is a bounded linear mapping from $W^{2,p}(B_R)$ into itself for any $p \in (1, \infty)$, $h \in L^p(B_R)$ and if, as in the proof of Theorem 9.11, $R \le \delta$ we must have $\|T\| \le \frac{1}{2}$. Therefore, by the contraction mapping principle (Theorem 5.1), (9.45) has a unique solution $v \in W^{2,p}(B_R)$ for any $p \in [1,r]$. Hence, $\eta u \in W^{2,r}(\Omega)$, and, since $x_0 \in \Omega$ is arbitrary, we obtain $u \in W^{2,r}_{\mathrm{loc}}(\Omega)$. If now $r = q$, we are done.

Otherwise, the desired interior regularity follows by using the Sobolev imbedding theorem and repeating the above argument. The case of local boundary regularity is handled similarly with $x_0 \in T$ and the ball $B_R(x_0)$ replaced by the half-ball $B_R^+(0)$ as in the proof of Theorem 9.13. $\square$

The uniqueness assertion of Theorem 9.15 follows from Lemma 9.16, for if the operator $L$ satisfies the hypotheses of Theorem 9.15 and the functions $u$, $v \in W^{2,p}(\Omega)$ satisfy $Lu = Lv$ in $\Omega$, $u - v \in W_0^{1,p}(\Omega)$, we have, by Lemma 9.16, $u - v \in W^{2,q}(\Omega) \cap W_0^{1,q}(\Omega)$ for all $1 < q < \infty$. Now using the uniqueness result, Theorem 9.5, and the Sobolev imbedding (Theorem 7.10), we conclude $u = v$. From the uniqueness, we can derive an apriori bound which extends Corollary 9.14.

\textbf{Lemma 9.17.} \textit{Let the operator $L$ satisfy the hypotheses of Theorem 9.15. Then there exists a constant $C$ (independent of $u$) such that}

\[
(9.46)\qquad \|u\|_{2,p;\Omega} \le C\|Lu\|_{p;\Omega}
\]

\textit{for all $u \in W^{2,p}(\Omega) \cap W_0^{1,p}(\Omega)$, $1 < p < \infty$.}

\textit{Proof.} We argue by contradiction. If (9.46) is not true, there must exist a sequence $\{v_m\} \subset W^{2,p}(\Omega) \cap W_0^{1,p}(\Omega)$ satisfying

\[
\|v_m\|_{p;\Omega} = 1;\qquad \|Lv_m\|_{p;\Omega} \to 0.
\]

By virtue of the apriori estimate (Theorem 9.13), the compactness of the imbedding $W_0^{1,p}(\Omega) \to L^p(\Omega)$, and the weak compactness of bounded sets in $W^{2,p}(\Omega)$ (Problem 5.5), there exists a subsequence, which we relabel as $\{v_m\}$, converging weakly to a function $v \in W^{2,p}(\Omega) \cap W_0^{1,p}(\Omega)$ satisfying $\|v\|_{p;\Omega} = 1$. Since

\[
\int_\Omega gD^\alpha v_m \to \int_\Omega gD^\alpha v
\]

for all $|\alpha| \le 2$ and $g \in L^{p/(p-1)}(\Omega)$, we must have

\[
\int_\Omega gLv = 0
\]


% PDF page 255
\source{gilbarg-trudinger:page-0255}

for all $g \in L^{p/(p-1)}(\Omega)$; hence $Lv = 0$ and $v = 0$ by the uniqueness assertion, which contradicts the condition $\|v\|_p = 1$. $\square$

We are now in a position to prove Theorem 9.15. First we observe that if the principal coefficients $a^{ij} \in C^0(\bar{\Omega})$ and $p \geq 2$, the result follows directly from Theorems 8.9 and 8.12, and Lemma 9.16. In the general case, we replace $u$ by $u - \varphi$, to reduce to zero boundary values, and approximate the coefficients $a^{ij}$ uniformly by a sequence $\{a_m^{ij}\} \subset C^0(\bar{\Omega})$ with, in the case $p < 2$, the function $f$ being approximated in $L^p(\Omega)$ by a sequence $\{f_m\} \subset L^2(\Omega)$. If $\{u_m\}$ denotes the sequence of solutions of the corresponding Dirichlet problems, we infer from Lemma 9.17 that the sequence $\{u_m\}$ is bounded in $W^{2,p}(\Omega)$. Consequently, by Problem 5.5 again, a subsequence converges weakly in $W^{2,p}(\Omega) \cap W^{1,p}_0(\Omega)$ to a function $u$ that (by an argument similar to that in Lemma 9.17) satisfies $Lu = f$ in $\Omega$. $\square$

Theorem 9.15 may also be derived from a Fredholm alternative, which follows from Corollary 9.14; (Problem 9.9). When $p > n/2$, we obtain an existence theorem for continuous boundary values.

\medskip

\noindent\textbf{Corollary 9.18.} \textit{Let $\Omega$ be a $C^{1,1}$ domain in $\mathbb{R}^n$, and let the operator $L$ be strictly elliptic in $\Omega$ with coefficients $a^{ij} \in C^0(\bar{\Omega})$, $b^i$, $c \in L^\infty$, $i, j = 1, \ldots, n$ and $c \leq 0$. Then, if $f \in L^p(\Omega)$, $p > n/2$, $\varphi \in C^0(\partial\Omega)$, the Dirichlet problem $Lu = f$ in $\Omega$, $u = \varphi$ on $\partial\Omega$, has a unique solution $u \in W^{2,p}_{\mathrm{loc}}(\Omega) \cap C^0(\bar{\Omega})$.}

\textit{Proof.} The uniqueness assertion follows from Theorem 9.5 and Lemma 9.16; (in fact, it is clearly valid for arbitrary $\Omega$ and $p > 1$). To get the existence, we let $\{\varphi_m\} \subset W^{2,p}(\Omega)$ converge uniformly to $\varphi$ on $\partial\Omega$, and let $u_m \in W^{2,p}(\Omega)$ be the solution of the Dirichlet problem $Lu_m = f$ in $\Omega$, $u_m = \varphi_m$ on $\partial\Omega$, known to exist by Theorem 9.15. The differences, $u_l - u_m$, clearly satisfy

\[
L(u_l - u_m) = 0 \quad \text{in } \Omega, \qquad u_l - u_m = \varphi_l - \varphi_m \quad \text{on } \partial\Omega.
\]

Hence by Theorems 9.1 and 9.11, and Lemma 9.16, we obtain that $\{u_m\}$ converges in $C^0(\bar{\Omega}) \cap W^{2,p}_{\mathrm{loc}}(\Omega)$ to a solution of the Dirichlet problem, $Lu = f$ in $\Omega$, $u = \varphi$ on $\partial\Omega$. $\square$

\medskip

By invoking barrier considerations similar to those in Chapter 6, we may extend Corollary 9.18 to allow for more general domains $\Omega$. We consider this type of result in Section 9.7 in conjunction with continuity estimates at the boundary.

To conclude this section we state a theorem concerning higher-order regularity that improves Theorems 6.17 and 6.19 for classical solutions. The proof may be effected by using difference quotients as in these theorems, or by an argument similar to that of Lemma 9.16. The details are left to the reader; (Problem 9.10).

\medskip

\noindent\textbf{Theorem 9.19.} \textit{Let $u$ be a $W^{2,p}_{\mathrm{loc}}(\Omega)$ solution of the elliptic equation $Lu = f$ in a domain $\Omega$, where the coefficients of $L$ belong to $C^{k-1,1}(\Omega)$, $(C^{k-1,\alpha}(\Omega))$, $f \in W^{k,q}_{\mathrm{loc}}(\Omega)$, $(C^{k-1,\alpha}(\Omega))$, with $1 < p, q < \infty, k \geq 1, 0 < \alpha < 1$. Then $u \in W^{k+2,q}_{\mathrm{loc}}(\Omega), (C^{k+1,\alpha}_{\mathrm{loc}}(\Omega)).$}


% PDF page 256
\source{gilbarg-trudinger:page-0256}

Furthermore, if $\Omega \in C^{k+1,\alpha}$, $(C^{k+1,\alpha})$, $L$ is strictly elliptic in $\Omega$ with coefficients in $C^{k-1,1}(\overline{\Omega})$, $(C^{k-1,\alpha}(\overline{\Omega}))$, and $f \in W^{k,q}(\Omega)$, $(C^{k-1,\alpha}(\overline{\Omega}))$, then $u \in W^{k+2,q}(\Omega)$, $(C^{k+1,\alpha}(\overline{\Omega}))$.

\begin{center}
\textbf{9.7. A Local Maximum Principle}
\end{center}

In this and the following sections, we focus attention on local pointwise estimates for operators in the general form (9.1) and derive results corresponding to those for divergence structure operators in Sections 8.6 through 8.10. We shall assume throughout the rest of this chapter that the operator $L$, as given by (9.1), is strictly elliptic with bounded coefficients in the domain $\Omega$, and accordingly, we fix constants $\gamma$ and $\nu$ such that

\[
(9.47)\qquad \frac{\Lambda}{\lambda} \le \gamma,\qquad \left(\frac{|b|}{\lambda}\right)^2,\ \frac{|c|}{\lambda} \le \nu
\]

in $\Omega$. We prove in this section the following analogue of the subsolution estimate, Theorem 8.17.

\textbf{Theorem 9.20.} \textit{Let $u \in W^{2,n}(\Omega)$ and suppose $Lu \ge f$, where $f \in L^n(\Omega)$. Then for any ball $B = B_{2R}(y) \subset \Omega$ and $p > 0$, we have}

\[
(9.48)\qquad \sup_{B_R(y)} u \le C\left\{\left(\frac{1}{|B|}\int_B (u^+)^p\right)^{1/p} + \frac{R}{\lambda}\|f\|_{L^n(B)}\right\}
\]

\textit{where $C = C(n,\gamma,\nu R^2,p)$.}

\textit{Proof.} Without loss of generality, we can assume that $B = B_1(0)$, the general case being recovered by means of the coordinate transformation $x \mapsto (x - y)/2R$. We also assume initially that $u \in C^2(\Omega) \cap W^{2,n}(\Omega)$. For $\beta \ge 1$, a cutoff function $\eta$ is defined by

\[
(9.49)\qquad \eta(x) = (1 - |x|^2)^\beta.
\]

By differentiation, we obtain

\[
D_i\eta = -2\beta x_i(1 - |x|^2)^{\beta - 1},
\]

\[
D_{ij}\eta = -2\beta\delta_{ij}(1 - |x|^2)^{\beta - 1} + 4\beta(\beta - 1)x_i x_j(1 - |x|^2)^{\beta - 2}.
\]

Setting $v = \eta u$, we then have

\[
a^{ij}D_{ij}v = \eta a^{ij}D_{ij}u + 2a^{ij}D_i\eta D_j u + ua^{ij}D_{ij}\eta
\]

\[
\ge \eta(f - b^iD_iu - cu) + 2a^{ij}D_i\eta D_j u + ua^{ij}D_{ij}\eta.
\]

% PDF page 257
\source{gilbarg-trudinger:page-0257}

\noindent Letting $\Gamma^{+} = \Gamma_i^{+}$ denote the upper contact set of $v$ in the ball $B$, we clearly have
$u > 0$ on $\Gamma^{+}$; and furthermore, using the concavity of $v$ on $\Gamma^{+}$, we can estimate
there

\[
|Du| = \frac{1}{\eta}\,|Dv - uD\eta|
\]
\[
\le \frac{1}{\eta}\,(|Dv| + u|D\eta|)
\]
\[
\le \frac{1}{\eta}\left(\frac{v}{1 - |x|} + u|D\eta|\right)
\]
\[
\le 2(1 + \beta)\eta^{-1/\beta}u.
\]

\noindent Thus, on $\Gamma^{+}$ we have the inequality

\[
-a^{ij}D_{ij}v \le \{(16\beta^2 + 2n\beta)\lambda\eta^{-2/\beta} + 2\beta|b|\eta^{-1/\beta} + c\}v + \eta f
\]
\[
\le C\lambda\eta^{-2/\beta}v + f,
\]

\noindent where $C = C(n,\beta,\gamma,\nu)$. Therefore, applying Lemma 9.3, we obtain for $\beta \ge 2$,

\[
\sup_B v \le C\left(\|\eta^{-2/\beta}v^{+}\|_{n;B} + \frac{1}{\lambda}\|f\|_{n;B}\right)
\]
\[
\le C\left\{(\sup_B v^{+})^{1-2/\beta}\|(u^{+})^{2/\beta}\|_{n;B} + \frac{1}{\lambda}\|f\|_{n;B}\right\}.
\]

\noindent Choosing $\beta = 2n/p$ (provided $p \le n$) and using Young's inequality (7.6) in the form

\[
(\sup_B v^{+})^{1-2/\beta} \le \epsilon \sup_B v^{+} + \epsilon^{1-\beta/2}
\]

\noindent for $\epsilon > 0$, we then get

\[
\sup_B v \le C\left\{\left(\int_B (u^{+})^{p}\right)^{1/p} + \frac{1}{\lambda}\|f\|_{n;B}\right\}
\]

\noindent and the estimate (9.48) follows. The extension to $u \in W^{2,n}(\Omega)$ follows directly by
approximation and is left to the reader. $\square$

\medskip

\noindent By replacement of $u$ with $-u$, Theorem 9.20 extends automatically to super-
solutions and solutions of the equation $Lu = f$.


% PDF page 258
\source{gilbarg-trudinger:page-0258}



\noindent\textbf{Corollary 9.21.}
Let $u\in W^{2,n}(\Omega)$ and suppose $Lu\le f$ $(=f)$, where $f\in L^n(\Omega)$.
Then, for any ball $B=B_{2R}(y)\subset\Omega$ and $p>0$, we have

\begin{equation}\tag{9.50}
\sup_{B_R(y)}(-u)\le C\left\{\left(\frac{1}{|B|}\int_B (u^-)^p\right)^{1/p}
+\frac{R}{\lambda}\,\|f\|_{L^n(B)}\right\}
\end{equation}

\noindent where $C=C(n,\gamma,\nu R^2,p)$.

\medskip

Note that when we take $p=1$ in (9.48), we obtain an extension of the mean value
inequality for non-negative subharmonic functions, namely,

\begin{equation}\tag{9.51}
u(y)\le \frac{C}{R^n}\int_{B_R(y)}u
\end{equation}

\noindent provided $Lu,u\ge 0$ in $B_R(y)$ and $C=C(n,\gamma,\nu R^2)$.
Theorem 9.20 also continues to hold under more general coefficient conditions;
(see Notes).

\bigskip

\noindent\textbf{9.8. \ Hölder and Harnack Estimates}

\medskip

We present in this section a treatment of the H\"older and Harnack estimates, of
Krylov and Safonov [KS 1, 2], that are the analogues for uniformly elliptic
operators in the general form of the De Giorgi, Nash and Moser estimates for
divergence form operators. Also important for our treatment of fully nonlinear
elliptic operators in Chapter 17 is the weak Harnack inequality for non-negative
supersolutions from which the H\"older and Harnack estimates are readily derived.

\medskip

\noindent\textbf{Theorem 9.22.}
Let $u\in W^{2,n}(\Omega)$ satisfy $Lu\le f$ in $\Omega$, where $f\in L^n(\Omega)$, and suppose
that u is non-negative in a ball $B=B_{2R}(y)\subset\Omega$. Then

\begin{equation}\tag{9.52}
\left(\frac{1}{|B_R|}\int_{B_R}u^p\right)^{1/p}\le C\left(\inf_{B_R}u+\frac{R}{\lambda}\,\|f\|_{L^n(B)}\right)
\end{equation}

\noindent where $p$ and $C$ are positive constants depending only on $n$, $\gamma$ and $\nu R^2$.

\medskip

\noindent\textit{Proof.} Again let us assume initially that $B=B_1(0)$ and also that $\lambda\equiv 1$, (by replacing $L$, $f$ by $L/\lambda$, $f/\lambda$). Setting

\begin{equation}\tag{9.53}
\bar u=u+\varepsilon+\|f\|_{n;B},\qquad
w=-\log \bar u,\qquad v=\eta w,\qquad g=f/\bar u,
\end{equation}


% PDF page 259
\source{gilbarg-trudinger:page-0259}



where $\varepsilon > 0$, and $\eta$ is given by (9.49), we then obtain, using Schwarz's inequality,
\[
-a^{ij}D_{ij}v = -\eta a^{ij}D_{ij}w - 2a^{ij}D_i\eta D_j w - wa^{ij}D_{ij}\eta
\]
\[
\le \eta\bigl(-a^{ij}D_i w D_j w + b^i D_i w + |c| + g\bigr)
- 2a^{ij}D_i\eta D_j w - wa^{ij}D_{ij}\eta
\]
\[
\le \frac{2}{\eta}a^{ij}D_i\eta D_j\eta - wa^{ij}D_{ij}\eta + \bigl(|b|^2 + |c| + g\bigr).
\]

\medskip

\noindent Next we calculate
\[
a^{ij}D_{ij}\eta = -2\beta a^{ii}(1 - |x|^2)^{\beta-1}
+ 4\beta(\beta - 1)a^{ij}x_i x_j(1 - |x|^2)^{\beta-2}
\]
\noindent so that $a^{ij}D_{ij}\eta \ge 0$ if
\[
2(\beta - 1)a^{ij}x_i x_j + a^{ii}|x|^2 \ge a^{ii}
\]
\noindent in particular if
\[
2\beta |x|^2 \ge n\Lambda.
\]

\medskip

\noindent Consequently if $0 < \alpha < 1$ and $\beta$ is chosen so that
\[
\beta \ge \frac{n\gamma}{2\alpha}
\]
\noindent we have $a^{ij}D_{ij}\eta \ge 0$ for all $|x| \ge \alpha$. Hence we obtain, on $B^+ = \{x\in B\,|\,w(x) > 0\}$,
the inequality
\[
-a^{ij}D_{ij}v \le 4\beta^2(1 - |x|^2)^{\beta-2}|x|^2
\]
\[
\qquad + v\chi(B_\alpha)\sup_{B_\alpha}\left(-\frac{a^{ij}D_{ij}\eta}{\eta}\right) + \bigl(|b|^2 + |c| + g\bigr)\eta
\]
\[
\le 4\beta^2 \Lambda + |b|^2 + |c| + g + \frac{2n\beta \Lambda}{1 - \alpha^2}v\chi(B_\alpha).
\]

\medskip

\noindent Therefore, applying Lemma 9.3, and noting that $\|g\|_{n;B} \le 1$, we get a bound for $v$,
\noindent namely,

\begin{equation}\tag{9.54}
\sup_B v \le C\bigl(1 + \|v^+\|_{n;B_\alpha}\bigr),
\end{equation}

\noindent where $C = C(n,\alpha,\gamma,\nu)$.

\medskip

\noindent To facilitate eventual application of the cube decomposition procedure in
Section 9.2 it is convenient at this point to switch from balls to cubes. For any point
$y \in \mathbb{R}^n$ and $R > 0$ we shall denote by $K_R(y)$ the open cube, parallel to the coordinate


% PDF page 260
\source{gilbarg-trudinger:page-0260}



\noindent axes, with center $y$ and side length $2R$. If $\alpha<\dfrac{1}{\sqrt n}$ we have $K_\alpha=K_\alpha(0)\subset B$ and
hence from (9.54),

\[
\sup_B v \le C(1+\|v^+\|_{n;K_\alpha})
\]
\[
\le C\bigl(1+|K_\alpha^+|^{1/n}\sup_B v^+\bigr),
\]

\noindent where $K_\alpha^+ = \{x\in K_\alpha\mid v>0\}$. Hence, if

\[
|K_\alpha^+|/|K_\alpha|\le \theta = [2(2\alpha)^n C]^{-1},
\]

\noindent then

\[
\sup_B v\le 2C,
\]

\noindent where $C=C(n,\alpha,\gamma,\nu)$ is the constant in (9.54). Let us now choose $\alpha=1/3n$ and
fix $\theta$ accordingly. Using the transformation $x\to\alpha(x-z)/r$, we therefore obtain for
any cube $K=K_r(z)$, such that $B_{3nr}(z)\subset B$ and

\begin{equation}\tag{9.55}
|K^+|\le \theta|K|,
\end{equation}

\noindent the estimate

\begin{equation}\tag{9.56}
\sup_{K_{3nr}(z)} w \le C(n,\gamma,\nu).
\end{equation}

\noindent The proof of Theorem 9.22 is now completed with the aid of the following
measure theoretic lemma.

\medskip

\noindent\textbf{Lemma 9.23.}
Let $K_0$ be a cube in $\mathbb{R}^n$, $w\in L^1(K_0)$; and, for $k\in\mathbb{R}$, set

\[
\Gamma_k=\{x\in K_0\mid w(x)\le k\}.
\]

\noindent Suppose there exist positive constants $\delta<1$ and $C$ such that

\begin{equation}\tag{9.57}
\sup_{K_0\cap K_{3r}(z)} (w-k)\le C,
\end{equation}

\noindent whenever $k$ and $K=K_r(z)\subset K_0$ satisfy

\begin{equation}\tag{9.58}
|\Gamma_k\cap K|\ge \delta|K|.
\end{equation}

\noindent Then it follows that, for all $k$,

\begin{equation}\tag{9.59}
\sup_{K_0}(w-k)\le C\left(1+\frac{\log(|\Gamma_k|/|K_0|)}{\log\delta}\right).
\end{equation}


% PDF page 261
\source{gilbarg-trudinger:page-0261}

\[
\textit{Proof.}\quad \text{We first show by induction that}
\]
\[
\sup_{K_0}(w-k)\le mC,
\]
for any natural number $m$ and $k\in\mathbb{R}$ satisfying $|\Gamma_k|\ge \delta^m|K_0|$. This assertion is clearly true by hypothesis for $m=1$. Suppose now it holds for some $m\in\mathbb{N}$, and that $|\Gamma_k|\ge \delta^{m+1}|K_0|$. Let $\widetilde{\Gamma}_k$ be defined by
\[
\widetilde{\Gamma}_k=\bigcup\{K_{3r}(z)\cap K_0\,|\,|K_r(z)\cap\Gamma_k|\ge \delta|K_r(z)|\}.
\]

By our cube decomposition procedure in Section 9.2, in particular inequality (9.20) with $t=\delta$, we obtain that either $\widetilde{\Gamma}_k=K_0$ or that
\[
|\widetilde{\Gamma}_k|\ge \frac{1}{\delta}|\Gamma_k|
\]
\[
\ge \delta^m|K_0|
\]
and hence, replacing $k$ by $k+C$, we obtain
\[
\sup_{K_0}(w-k)\le (m+1)C,
\]
which guarantees the validity of the above assertion for $m+1$. The estimate (9.59) now follows by the appropriate choice of $m$. $\square$

To apply Lemma 9.23, we take $\delta=1-\theta$, $K_0=K_\alpha(0)$, $\alpha=1/3n$ and note that the estimate (9.56) still holds when $w$ is replaced by $w-k$. Let
\[
\mu_t=\bigl|\{x\in K_0\,|\,\bar{u}(x)>t\}\bigr|
\]
denote the distribution function of $\bar{u}$ in $K_0$ to obtain from (9.53) and (9.59) with $t=e^{-k}$, the estimate
\[
(9.60)\qquad \mu_t\le C\bigl(\inf_{K_0}\bar{u}/t\bigr)^{\kappa},\qquad t>0,
\]
where $C$ and $\kappa$ are positive constants depending only on $n$, $\gamma$ and $\nu$. Replacing the cube $K_0$ by the inscribed ball $B_\alpha(0)$, $\alpha=1/3n$ and using Lemma 9.7, we then obtain
\[
(9.61)\qquad \int_{B_\alpha}(\bar{u})^p\le C\bigl(\inf_{B_\alpha}\bar{u}\bigr)^p,
\]
for $p<\kappa$, say $p=\kappa/2$. The weak Harnack inequality in the form (9.52) then follows by letting $\epsilon\to 0$, using a covering argument to extend (9.61) to arbitrary $\alpha<1$, (in particular $\alpha=1/2$), and finally invoking the coordinate transformation $x\to (x-y)/2R$. $\square$

% PDF page 262
\source{gilbarg-trudinger:page-0262}

By adapting the proof of the Hölder estimate for divergence structure operators
(Theorem 8.22), (with only a minor modification to offset not having $p=1$ in the
weak Harnack inequality (9.52)), we conclude from Theorem 9.22 the following
Hölder estimate for operators in general form.

\noindent\textbf{Corollary 9.24.} \quad Let $u \in W^{2,n}(\Omega)$ satisfy the equation $Lu = f$ in $\Omega$. Then, for any ball
$B_0 = B_{R_0}(y) \subset \Omega$ and $R \leq R_0$, we have

\begin{equation}
\osc_{B_R(y)} u \leq C \left(\frac{R}{R_0}\right)^{\alpha} \left(\osc_{B_0} u + kR_0\right),
\tag{9.62}
\end{equation}

where $C = C(n,\gamma,\nu R_0^2)$, $\alpha = \alpha(n,\gamma,\nu R_0^2)$ are positive constants and
$k = \|f - cu\|_{n;B_0}$.

Also by combining Theorem 9.22 with the subsolution estimate (Theorem 9.20),
we obtain the full Harnack inequality.

\noindent\textbf{Corollary 9.25.} \quad Let $u \in W^{2,n}(\Omega)$ satisfy $Lu = 0$, $u \geq 0$ in $\Omega$. Then for any ball
$B_{2R}(y) \subset \Omega$, we have

\begin{equation}
\sup_{B_R(y)} u \leq C \inf_{B_R(y)} u,
\tag{9.63}
\end{equation}

where $C = C(n,\gamma,\nu R^2)$.

\section*{9.9. Local Estimates at the Boundary}

The local maximum principle (Theorem 9.20) may be extended to balls inter-
secting the boundary $\partial\Omega$ as follows.

\noindent\textbf{Theorem 9.26.} \quad Let $u \in W^{2,n}(\Omega) \cap C^0(\overline{\Omega})$ satisfy $Lu \geq f$ in $\Omega$, $u \leq 0$ on $B \cap \partial\Omega$
where $f \in L^n(\Omega)$ and $B = B_{2R}(y)$ is a ball in $\R^n$. Then, for any $p > 0$, we have

\begin{equation}
\sup_{\Omega \cap B_R(y)} u \leq C \left\{\left(\frac{1}{|B|}\int_{B \cap \Omega}(u^+)^p\right)^{1/p}
+ \frac{R}{\lambda}\|f\|_{L^n(B \cap \Omega)}\right\},
\tag{9.64}
\end{equation}

where $C = C(n,\gamma,\nu R^2,p)$.

Proof. \quad It suffices to establish the estimate (9.64) for $u \in C^2(\Omega) \cap C^0(\overline{\Omega})$ satisfying
$u \leq 0$ on $B \cap \partial\Omega$. We extend $u$ to the whole of the ball $B$ by setting $u = 0$ in $B - \Omega$.
Although $u$ then does not necessarily belong to $C^2(B)$, the argument of Theorem
9.20 may still be applied since $\Gamma^+$, the upper contact set of the function $v$, will lie
in $B \cap \Omega$, and $v \in C^2(\Gamma^+)$ is sufficient for the application of Lemma 9.3. \hfill $\square$

It is worth recording here the general observation, used in the proof of Theorem
9.26, that if $u$ fulfills the hypotheses of Lemma 9.3, the estimate (9.11) continues to


% PDF page 263
\source{gilbarg-trudinger:page-0263}

hold when $\Gamma^{+}$ is replaced by the upper contact set of $u$ with respect to any larger domain $\widetilde{\Omega}$, (to which $u$ is extended by setting $u=0$ in $\widetilde{\Omega}-\Omega$), and with $d=\operatorname{diam}\widetilde{\Omega}$.
The weak Harnack inequality (Theorem 9.22) admits the following extension to the boundary.

\medskip

\noindent\textbf{Theorem 9.27.} \textit{Let $u\in W^{2,n}(\Omega)$ satisfy $Lu\le f$ in $\Omega$, $u\ge 0$ in $B\cap\Omega$, where $B=B_{2R}(y)$ is a ball in $\mathbb{R}^n$. Set}
\[
m=\inf_{B\cap\Omega}u
\]
\textit{and}
\[
u_m^{-}(x)=
\begin{cases}
\inf\{u(x),m\} & \text{for } x\in B\cap\Omega,\\
m & \text{for } x\in B-\Omega.
\end{cases}
\]
\textit{Then}
\begin{equation}
\left(\frac{1}{|B_R|}\int_{B_R}(u_m^{-})^p\right)^{1/p}
\le
C\left(\inf_{\Omega\cap B_R}u+\frac{R}{\lambda}\|f\|_{L^n(B\cap\Omega)}\right),
\tag{9.65}
\end{equation}
\textit{where $p$ and $C$ are positive constants depending only on $n,\gamma,$ and $\nu R^2$. If we assume only $u\in W^{2,n}_{\mathrm{loc}}(\Omega)$, then (9.65) holds with}
\[
m=\lim_{x\to B\cap\partial\Omega}\inf u.
\]

\medskip

\noindent\textit{Proof.} We adapt the proof of Theorem 9.22 by replacing $u$ with $u_m^{-}$. The estimate (9.56) then follows when the function $w$ is replaced by $w-k$ for $k\ge -\log m$, and we thus infer the estimate (9.60) for $0<t\le m$. But $\mu_t=0$, if $t>m$, hence (9.65) follows as before. The final assertion of Theorem 9.27 is a consequence of the remark following Theorem 9.1. $\square$

\medskip

Global and boundary modulus of continuity estimates follow as consequences of Theorem 9.27. Corresponding to the divergence structure results (Theorems 8.27 and 8.29), we obtain the following estimates.

\medskip

\noindent\textbf{Corollary 9.28.} \textit{Let $u\in W^{2,n}_{\mathrm{loc}}(\Omega)$ satisfy $Lu=f$ in $\Omega$ where $f\in L^n(\Omega)$, and suppose that $\Omega$ satisfies an exterior cone condition at a point $y\in\partial\Omega$. Then, for any $0<R<R_0$ and ball $B_0=B_{R_0}(y)$, we have}
\begin{equation}
\osc_{\Omega\cap B_R}u
\le
C\left\{\left(\frac{R}{R_0}\right)^{\alpha}\left(\osc_{\Omega\cap B_0}u+kR_0\right)+\sigma\sqrt{RR_0}\right\},
\tag{9.66}
\end{equation}
\textit{where $C=C(n,\gamma,\nu R_0^2,V_y)$, $\alpha=\alpha(n,\gamma,\nu R_0^2,V_y)$ are positive constants, $V_y$ is the exterior cone at $y$ and}
\[
\sigma(r)=\osc_{\Omega\cap B_r}u=\limsup_{x\to\partial\Omega\cap B_r}u-\liminf_{x\to\partial\Omega\cap B_r}u
\]
\textit{for $0<r\le R_0$.}


% PDF page 264
\source{gilbarg-trudinger:page-0264}
\providecommand{\loc}{\mathrm{loc}}

\begin{quote}
\textbf{Corollary 9.29.} \emph{Let $u\in W^{2,n}(\Omega)\cap C^{0}(\overline{\Omega})$ satisfy $Lu=f$ in $\Omega$, $u=\varphi$ on $\partial\Omega$ where $f\in L^{n}(\Omega)$, $\varphi\in C^{\beta}(\overline{\Omega})$ for some $\beta>0$, and suppose that $\partial\Omega$ satisfies a uniform exterior cone condition. Then $u\in C^{\alpha}(\overline{\Omega})$ and}

\[
(9.67)\qquad |u|_{\alpha;\Omega}\le C,
\]

\emph{where $\alpha$ and $C$ are positive constants depending on $n$, $\gamma$, $\nu$, $\beta$, $\Omega$, $|\varphi|_{s;\Omega}$ and $|u|_{0;\Omega}$.}

We note here that modulus of continuity estimates at the boundary also arise from barrier constructions as in Section 6.3. However, Theorem 9.28 may be used in place of barrier arguments to solve the Dirichlet problem by means of the Perron process. To see this, let us suppose that the operator $L$ satisfies the hypotheses of Theorem 6.11; and let $u\in C^{2}(\Omega)$ be the Perron solution of the Dirichlet problem, $Lu=f$, $u=\varphi$ on $\partial\Omega$, whose existence is asserted by Theorem 6.11. Corollary 9.28 then provides an estimate for the modulus of continuity of $u$ at a point $y\in \partial\Omega$, where $\Omega$ satisfies an exterior cone condition, in terms of the modulus of continuity of the function $\varphi$ at $y$. If $\Omega$ satisfies an exterior cone condition everywhere on $\partial\Omega$, we then conclude $u\in C^{0}(\overline{\Omega})$ and $u=\varphi$ on $\partial\Omega$ thus solving the above Dirichlet problem.

Consequently, in the existence theorem (Theorem 6.13), we may replace the exterior sphere condition by an exterior cone condition that would be satisfied, for example, by Lipschitz domains; (see also Problem 6.3). Utilizing our results in Section 9.5, in particular Corollary 9.18, we may further extend Theorem 6.13 to cover continuous coefficients.

\medskip

\textbf{Theorem 9.30.} \emph{Let $L$ be strictly elliptic in a bounded domain $\Omega$ with coefficients $a^{ij}\in C^{0}(\Omega)\cap L^{\infty}(\Omega)$, $b^{i},c\in L^{\infty}(\Omega)$ and $c\le 0$, and suppose that $\Omega$ satisfies an exterior cone condition at every boundary point. Then if $f\in L^{p}(\Omega)$, $p\ge n$, the Dirichlet problem, $Lu=f$ in $\Omega$, $u=\varphi$ on $\partial\Omega$, has a unique solution $u\in W^{2,p}_{\loc}(\Omega)\cap C^{0}(\overline{\Omega})$.}

\emph{Proof.} To complete the proof of Theorem 9.30, it suffices, in view of the remarks preceding the theorem, to establish the existence of an analogue in $W^{2,p}_{\loc}(\Omega)$ of a Perron solution under the hypotheses of the theorem. Again, we may imitate the Perron method for subharmonic functions in Chapter 2 making crucial use of the strong maximum principle (Theorem 9.6), the solvability of the Dirichlet problem in balls with continuous boundary data (Corollary 9.18), and the interior estimates (Theorem 9.11) coupled with the weak relative compactness of bounded sets in $W^{2,p}(\Omega')$, $\Omega'\Subset\Omega$. The details are left to the reader. \qed
\end{quote}

\textbf{Boundary Hölder Estimates for the Gradient}

An interesting and useful Hölder estimate for the traces of the gradients of solutions on the boundary may also be deduced from the interior Harnack (or weak Harnack) inequalities. This result was established by Krylov [KV 5] in connection with its application to the theory of fully nonlinear equations, which we describe in Section 17.8. For these purposes it suffices to restrict to a flat boundary portion

% PDF page 265
\source{gilbarg-trudinger:page-0265}

where the solution is assumed to vanish and to consider operators of the form

\[
Lu = a^{ij} D_{ij}u
\]

satisfying the uniform ellipticity condition (9.47). More general results are readily formulated.

\medskip

\noindent\textbf{Theorem 9.31.} \textit{Let $u \in W^{2,n}_{\mathrm{loc}}(B^+) \cap C^0(\overline{B^+})$ satisfy the equation $Lu = f$ in the half ball $B^+ = B_{R_0}(0) \cap \R^n_+$ with $f \in L^\infty(B^+)$ and $u = 0$ on $T = B_{R_0} \cap \partial \R^n_+$. Then for any $R \leq R_0$, we have}

\begin{equation}\tag{9.68}
\osc_{B_R^+} \frac{u}{x_n} \leq C \left(\frac{R}{R_0}\right)^{\alpha}
\left(\osc_{B^+} \frac{u}{x_n} + R_0 \sup_{B^+} \left|\frac{f}{\lambda}\right|\right)
\end{equation}

\noindent\textit{where $\alpha$ and $C$ are positive constants depending only on $n$ and $\gamma$.}

\medskip

\noindent\textit{Proof.} We proceed under the assumption that the function $v = u/x_n$ is bounded in $B^+$; (local boundedness at least is guaranteed by the barrier considerations in Section 6.3). Assuming initially that $u \geq 0$ in $B^+$, we first prove the following assertion: There exists $\delta = \delta(n,\gamma) > 0$ such that

\begin{equation}\tag{9.69}
\inf_{\substack{|x'|<R\\ x_n=\delta R}} v \leq 2\left(\inf_{B_{R/2,\delta}} v + \frac{R}{\lambda}\sup_{B^+}|f|\right)
\end{equation}

\noindent for any $R \leq R_0$, where

\[
B_{R,\delta}=\{x \mid |x'|<R,\ 0<x_n<\delta R\}.
\]

To prove (9.69) it is convenient to normalize so that $\lambda = R = 1$ and
\[
\inf_{|x'|<R} v(x',\delta R)=1.
\]
We consider in $B_{1,\delta}$ the barrier function

\[
w(x)=\left(1-|x'|^2 + (1+\sup|f|)\frac{x_n-\delta}{\sqrt{\delta}}\right)x_n.
\]

By straightforward computation, we obtain $Lw \geq f$ for sufficiently small $\delta = \delta(n,\gamma)$ and $w \leq u$ on $\partial B_{1,\delta}$, so by the maximum principle, Theorem 9.1, we have $w \leq u$ in $B_{1,\delta}$. Thus on $B_{1/2,\delta}$,

\[
v \geq 1 - |x'|^2 + (1+\sup|f|)\frac{x_n-\delta}{\sqrt{\delta}} \geq \frac{1}{2} - \sup|f|,
\]

again for sufficiently small $\delta$. Removal of the normalization yields (9.69) as required. We now define

\[
B^*_{R/2,\delta}=\{x \mid |x'|<R,\ \delta R/2 < x_n < 3\delta R/2\},
\]

noting that

\[
\frac{2u}{3\delta R}\leq v \leq \frac{2u}{\delta R}.
\]


% PDF page 266
\source{gilbarg-trudinger:page-0266}

in $B^*_{R/2,\delta}$. Consequently by the Harnack inequality, Corollary 9.25, we obtain

\begin{equation}
\tag{9.70}
\sup_{B^*_{R/2,\delta}} v \leq C\left(\inf_{B^*_{R/2,\delta}} v + R \sup |f/\lambda|\right)
\end{equation}

\[
\leq C\left(\inf_{\substack{|x'|<R,\\ x_n=\delta R}} v + R \sup |f/\lambda|\right)
\]

\[
\leq C\left(\inf_{B_{R,\delta}} v + R \sup |f/\lambda|\right), \qquad \text{by (9.69).}
\]

We now drop the assumption $u \geq 0$ and set $M = \sup_{B_{2R,\delta}} v$, $m = \inf_{B_{2R,\delta}} v$. Applying (9.70)
to the functions $M - v$, $v - m$ and adding the resulting inequalities, we get the
standard oscillation estimate

\[
\osc_{B^*_{R/2,\delta}} v \leq \sigma\left(\osc_{B_{R,\delta}} v + C R \sup |f/\lambda|\right)
\]

where $C > 0$ and $\sigma < 1$ depend only on $n$ and $\gamma$. From this (9.68) follows by virtue
of Lemma 8.23. $\square$
Theorem 9.31 shows in fact that the gradient $Du$ exists on $T$ and is H\"older
continuous there, satisfying the estimate

\begin{equation}
\tag{9.71}
\osc_{|x'|<R} Du(x',0) \leq C\left(\frac{R}{R_0}\right)^{\alpha}
\left(\osc_{B^+} \frac{u}{x_n} + \frac{R}{\lambda} \sup_{B^+} |f|\right)
\end{equation}

Moreover the term $\osc_{B^+} \frac{u}{x_n}$ on the right hand side of (9.68) or (9.71) can be replaced
by either $\osc_{B^+} u/R_0$ or $\sup_{B^+} |Du|$. Global $C^{1,\alpha}$ estimates can be deduced from Theo-
rem 9.31 on combination with appropriate interior estimates which typically hold
for nonlinear equations; (see Problem 13.1 and Section 17.8).

\bigskip

\noindent\textbf{Notes}

\medskip

The maximum and uniqueness principles for strong solutions, as formulated in
Theorems 9.1, 9.5 and 9.6, are due to Aleksandrov [AL 2], [AL 3], although the
essential case covered by Lemma 9.3 appears in Bakelman [BA 3]. A detailed
analysis of the form of the constant $C$ in the estimate (9.4) is carried out in the papers
[AL 4, 5]. In all these results it is impossible to replace the $L^n$ norms by $L^p$ norms
for $p < n$ [AL 6]. Indeed the example (8.22) shows that the uniqueness result
(Theorem 9.5) no longer holds if we only assume the functions $u, v \in W^{2,p}_{\mathrm{loc}}(\Omega) \cap$
$C^0(\overline{\Omega})$, $p < n$. Different versions of the maximum principle (Theorem 9.1) were
discovered by Bony [BY] and Pucci [PU 3].

The Calderon-Zygmund inequality was discovered by Calderon and Zygmund
[CZ], and we have largely followed their original proof, as expounded by Stein
[SN], making use of the cube decomposition procedure (which extends the one
dimensional case due to Riesz [RZ]) and the Marcinkiewicz interpolation theorem
[MZ]. Our proof differs from those in [CZ] and [SN] in that we do not use the


% PDF page 267
\source{gilbarg-trudinger:page-0267}

Fourier transform to get the $L^2$ estimate. The operator $T$ in the proof of Theorem
9.9 is a special case of a singular integral operator, and it was this class of operators
that was the main object of study in [CZ] and [SN]. An alternative proof of the
Calderon-Zygmund inequality is presented in the monographs [BS], [MY 5]. A
further proof, also based on interpolation, uses the space of functions of bounded
mean oscillation, BMO; see [CS], [FS].

The $L^p$ estimates for second-order elliptic equations, presented in Section 9.4,
were derived by Koselev [KO] and Greco [GC] and extended to higher-order
equations and systems by various authors, including Slobodeckii [SL] Browder
[BW 3] and Agmon, Douglis and Nirenberg [ADN 1, 2]. The existence theorem
(Theorem 9.15) appears in Chicco [CI 4, 5], although derived differently there; and
a further approach to the Dirichlet problem is given by P.\ L.\ Lions [LP 1]. The
regularity argument in Lemma 9.16 is taken from Morrey [MY 5], where the $L^p$
theory is also treated.

The pointwise estimates of Sections 9.6, 9.7 and 9.8 stem from the fundamental
work of Krylov and Safonov, (see [KS 1, 2], [SF]) where the H\"older and Harnack
estimates of Corollaries 9.24, 9.25 are established, for $c \le 0$. In fact, the more
general situation of parabolic equations is treated in [KS 1, 2]. Our presentation in
Section 9.7, adapted from [TR 12], carries over their basic ideas. The local maximum
principle (Theorem 9.20) was also proved in [TR 12], under more general coefficient
conditions, namely $A/\Dop^*$, $b/\Dop^* \in L^q(\Omega)$, $q > n$, $c/\Dop^*$, $f/\Dop^* \in L^n(\Omega)$. The estimates
of Section 9.7 may similarly be extended to allow $b/\lambda \in L^{2n}(\Omega)$, $c/\lambda$, $f/\lambda \in L^n(\Omega)$,
although the condition of uniform ellipticity seems essential for the proofs in this
case. Extensions to quasilinear equations are treated in [TR 12], [LU 7] and [MV 1];
(see also Chapter 15).

We have also included in the present edition a proof of the boundary H\"older
gradient estimate of Krylov [KV 5], utilizing a simplification of Caffarelli. Krylov's
original proof was presented in the form of a problem in the English second
edition.

\bigskip

\noindent\textbf{Problems}

\medskip

\noindent\textbf{9.1.} \enspace Prove that in the estimates (9.8) and (9.11) we can replace $d$ by $d/2$ and also
by considering a spherical cone show that the resulting estimates are sharp; (see
[AL 4, 5]).

\medskip

\noindent\textbf{9.2.} \enspace By explicitly integrating the function $g$ and optimizing the choice of $\mu$ in the
proof of Theorem 9.1, deduce an improvement of the estimate (9.14); (see [AL 4, 5]).

\medskip

\noindent\textbf{9.3.} \enspace Derive the estimate (9.11) in the form

\begin{equation}
\tag{9.68}
\sup_{\Omega} u \le \sup_{\partial \Omega} u^+ + C(n) |\widehat{\Omega}|^{1/n}
\left\|\frac{a^{ij}D_{ij}u}{\Dop^*}\right\|_{L^n(\Gamma^+)} ,
\end{equation}

where $\widehat{\Omega}$ denotes the convex hull of $\Omega$; (see [TA 5] for a sharp version of (9.68) when
$n = 2$).


% PDF page 268
\source{gilbarg-trudinger:page-0268}

\begin{quote}
\textbf{9.4.} Prove the following more general version of the Marcinkiewicz interpolation theorem:

\textbf{Theorem 9.31.} \emph{Let $T$ be a linear mapping from $L^q(\Omega)\cap L^r(\Omega)$ into $L^{\bar q}(\Omega)\cap L^{\bar r}(\Omega)$, $1\le q<r<\infty$, $1\le \bar q\le \bar r<\infty$, and suppose there are constants $T_1$ and $T_2$ such that}

\[
\mu_{Tf}(t)\le \left(\frac{T_1\lVert f\rVert_q}{t}\right)^{\bar q},
\qquad
\mu_{Tf}(t)\le \left(\frac{T_2\lVert f\rVert_r}{t}\right)^{\bar r}
\tag{9.69}
\]

\emph{for all $f\in L^q(\Omega)\cap L^r(\Omega)$ and $t>0$. Then $T$ extends as a bounded mapping from $L^p(\Omega)$ to $L^{\bar p}(\Omega)$ for any $p,\bar p$ satisfying}

\[
1/p=\sigma/q+(1-\sigma)/r,\qquad 1/\bar p=\sigma/\bar q+(1-\sigma)/\bar r
\]

\emph{for some $\sigma\in(0,1)$.}

\emph{Examine also the case when $r$ or $\bar r=\infty$.}

\textbf{9.5.} With the notation of Section 7.8, use the general Marcinkiewicz interpolation theorem to show that the potential operator $V_\mu$ maps $L^p(\Omega)$ continuously into $L^q(\Omega)$ for $p>1$ and $\delta=\mu$.

\textbf{9.6.} Using Lemma 9.12, show that for a $C^{1,1}$ domain $\Omega$ the subspace

\[
\{u\in C^2(\bar\Omega)\mid u=0\quad \text{on }\partial\Omega\}
\]

is dense in $W^{2,p}(\Omega)\cap W^{1,p}_0(\Omega)$, $1<p<\infty$.

\textbf{9.7.} Deduce Theorem 9.15 directly by approximation based on either Theorems 6.14 or 8.14.

\textbf{9.8.} Deduce Theorem 9.15 for the special case of Poisson's equation from the Riesz representation theorem, and use the method of continuity (Theorem 5.2) to get the full result.

\textbf{9.9.} Starting from Corollary 9.14, establish a Fredholm alternative for operators of the form (9.1) in the Sobolev spaces $W^{2,p}(\Omega)$, $1<p<\infty$; and again deduce Theorem 9.15.

\textbf{9.10.} Prove Theorem 9.19.

\textbf{9.11.} Suppose the operator $L$ satisfies the hypotheses of Theorem 9.22 in an annular region, $A=B_R(y)-\overline{B}_\rho(y)\subset \mathbb{R}^n$. If $u\in W^{2,n}(A)$ satisfies $Lu\le f$, $u\ge 0$ in $A$, where $f\in L^n(A)$, prove that for any $\rho<r<R$

\[
\tag{9.70}
\inf_{B_r-B_\rho}u\ge \kappa\bigl(\inf_{\partial B_\rho}u-R\lVert f\rVert_{n;A}\bigr),
\]
\end{quote}

% PDF page 269
\source{gilbarg-trudinger:page-0269}

where $\kappa$ is a positive constant depending only on $n$, $\rho/R$, $r/R$, $\gamma$ and $\nu R^2$. Use this result to deduce Theorem 9.22 from Theorem 9.20.

\medskip

\noindent\textbf{Problem 9.12.} By considering the operator

\[
Lu = \Lambda \Delta u + (\lambda - \Lambda) \frac{x_i x_j}{|x|^2} D_{ij}u
\]

\noindent and the function

\[
u(x) = |x|^{1-(n-1)\gamma}, \qquad \gamma = \Lambda/\lambda
\]

\noindent suitably redefined near $0$, show that the exponent $p$ in the weak Harnack inequality Theorem 9.22 must satisfy

\[
p < \frac{n}{(n-1)\gamma - 1}.
\]

% PDF page 270
\source{gilbarg-trudinger:page-0270}

% Part II
% Quasilinear Equations

% PDF page 271
\source{gilbarg-trudinger:page-0271}

\chapter{Maximum and Comparison Principles}

The purpose of this chapter is to provide various maximum and comparison principles for quasilinear equations which extend corresponding results in Chapter 3. We consider second order, quasilinear operators $Q$ of the form

\begin{equation}
Qu=a^{ij}(x,u,Du)D_{ij}u+b(x,u,Du), \qquad a^{ij}=a^{ji},
\tag{10.1}
\end{equation}

where $x=(x_1,\dots,x_n)$ is contained in a domain $\Omega$ of $\mathbb{R}^n$, $n\ge 2$, and, unless otherwise stated, the function $u$ belongs to $C^2(\Omega)$. The coefficients of $Q$, namely the functions $a^{ij}(x,z,p)$, $i,j=1,\dots,n$, $b(x,z,p)$ are assumed to be defined for all values of $(x,z,p)$ in the set $\Omega\times\mathbb{R}\times\mathbb{R}^n$. Two operators of the form (10.1) will be called \emph{equivalent} if one is a multiple of the other by a fixed positive function in $\Omega\times\mathbb{R}\times\mathbb{R}^n$. Equations $Qu=0$ corresponding to equivalent operators $Q$ will also be called equivalent.

We adopt the following definitions:

Let $\mathcal{U}$ be a subset of $\Omega\times\mathbb{R}\times\mathbb{R}^n$. Then $Q$ is \emph{elliptic} in $\mathcal{U}$ if the coefficient matrix $[a^{ij}(x,z,p)]$ is positive for all $(x,z,p)\in\mathcal{U}$. If $\lambda(x,z,p)$, $\Lambda(x,z,p)$ denote respectively the minimum and maximum eigenvalues of $[a^{ij}(x,z,p)]$, this means that

\begin{equation}
0<\lambda(x,z,p)|\xi|^2\le a^{ij}(x,z,p)\xi_i\xi_j\le \Lambda(x,z,p)|\xi|^2
\tag{10.2}
\end{equation}

for all $\xi=(\xi_1,\dots,\xi_n)\in\mathbb{R}^n-\{0\}$ and for all $(x,z,p)\in\mathcal{U}$. If, further, $\Lambda/\lambda$ is bounded in $\mathcal{U}$, we shall call $Q$ \emph{uniformly elliptic} in $\mathcal{U}$. If $Q$ is elliptic (uniformly elliptic) in the whole set $\Omega\times\mathbb{R}\times\mathbb{R}^n$, then we shall simply say that $Q$ is elliptic (uniformly elliptic) in $\Omega$. If $u\in C^1(\Omega)$ and the matrix $[a^{ij}(x,u(x),Du(x))]$ is positive for all $x\in\Omega$, we shall say that $Q$ is elliptic with respect to $u$. We also define a scalar function $\mathcal{E}$, which will prove to be quite important, by

\begin{equation}
\mathcal{E}(x,z,p)=a^{ij}(x,z,p)p_ip_j.
\tag{10.3}
\end{equation}

If $Q$ is elliptic in $\mathcal{U}$, we have by (10.2)

\begin{equation}
0<\lambda(x,z,p)|p|^2\le \mathcal{E}(x,z,p)\le \Lambda(x,z,p)|p|^2
\tag{10.4}
\end{equation}

for all $(x,z,p)\in\mathcal{U}$.
